% Options for packages loaded elsewhere
% Options for packages loaded elsewhere
\PassOptionsToPackage{unicode}{hyperref}
\PassOptionsToPackage{hyphens}{url}
\PassOptionsToPackage{dvipsnames,svgnames,x11names}{xcolor}
%
\documentclass[
  letterpaper,
  DIV=11,
  numbers=noendperiod,
  oneside]{scrartcl}
\usepackage{xcolor}
\usepackage[left=1in,marginparwidth=2.0666666666667in,textwidth=4.1333333333333in,marginparsep=0.3in]{geometry}
\usepackage{amsmath,amssymb}
\setcounter{secnumdepth}{-\maxdimen} % remove section numbering
\usepackage{iftex}
\ifPDFTeX
  \usepackage[T1]{fontenc}
  \usepackage[utf8]{inputenc}
  \usepackage{textcomp} % provide euro and other symbols
\else % if luatex or xetex
  \usepackage{unicode-math} % this also loads fontspec
  \defaultfontfeatures{Scale=MatchLowercase}
  \defaultfontfeatures[\rmfamily]{Ligatures=TeX,Scale=1}
\fi
\usepackage{lmodern}
\ifPDFTeX\else
  % xetex/luatex font selection
\fi
% Use upquote if available, for straight quotes in verbatim environments
\IfFileExists{upquote.sty}{\usepackage{upquote}}{}
\IfFileExists{microtype.sty}{% use microtype if available
  \usepackage[]{microtype}
  \UseMicrotypeSet[protrusion]{basicmath} % disable protrusion for tt fonts
}{}
\makeatletter
\@ifundefined{KOMAClassName}{% if non-KOMA class
  \IfFileExists{parskip.sty}{%
    \usepackage{parskip}
  }{% else
    \setlength{\parindent}{0pt}
    \setlength{\parskip}{6pt plus 2pt minus 1pt}}
}{% if KOMA class
  \KOMAoptions{parskip=half}}
\makeatother
% Make \paragraph and \subparagraph free-standing
\makeatletter
\ifx\paragraph\undefined\else
  \let\oldparagraph\paragraph
  \renewcommand{\paragraph}{
    \@ifstar
      \xxxParagraphStar
      \xxxParagraphNoStar
  }
  \newcommand{\xxxParagraphStar}[1]{\oldparagraph*{#1}\mbox{}}
  \newcommand{\xxxParagraphNoStar}[1]{\oldparagraph{#1}\mbox{}}
\fi
\ifx\subparagraph\undefined\else
  \let\oldsubparagraph\subparagraph
  \renewcommand{\subparagraph}{
    \@ifstar
      \xxxSubParagraphStar
      \xxxSubParagraphNoStar
  }
  \newcommand{\xxxSubParagraphStar}[1]{\oldsubparagraph*{#1}\mbox{}}
  \newcommand{\xxxSubParagraphNoStar}[1]{\oldsubparagraph{#1}\mbox{}}
\fi
\makeatother

\usepackage{color}
\usepackage{fancyvrb}
\newcommand{\VerbBar}{|}
\newcommand{\VERB}{\Verb[commandchars=\\\{\}]}
\DefineVerbatimEnvironment{Highlighting}{Verbatim}{commandchars=\\\{\}}
% Add ',fontsize=\small' for more characters per line
\usepackage{framed}
\definecolor{shadecolor}{RGB}{241,243,245}
\newenvironment{Shaded}{\begin{snugshade}}{\end{snugshade}}
\newcommand{\AlertTok}[1]{\textcolor[rgb]{0.68,0.00,0.00}{#1}}
\newcommand{\AnnotationTok}[1]{\textcolor[rgb]{0.37,0.37,0.37}{#1}}
\newcommand{\AttributeTok}[1]{\textcolor[rgb]{0.40,0.45,0.13}{#1}}
\newcommand{\BaseNTok}[1]{\textcolor[rgb]{0.68,0.00,0.00}{#1}}
\newcommand{\BuiltInTok}[1]{\textcolor[rgb]{0.00,0.23,0.31}{#1}}
\newcommand{\CharTok}[1]{\textcolor[rgb]{0.13,0.47,0.30}{#1}}
\newcommand{\CommentTok}[1]{\textcolor[rgb]{0.37,0.37,0.37}{#1}}
\newcommand{\CommentVarTok}[1]{\textcolor[rgb]{0.37,0.37,0.37}{\textit{#1}}}
\newcommand{\ConstantTok}[1]{\textcolor[rgb]{0.56,0.35,0.01}{#1}}
\newcommand{\ControlFlowTok}[1]{\textcolor[rgb]{0.00,0.23,0.31}{\textbf{#1}}}
\newcommand{\DataTypeTok}[1]{\textcolor[rgb]{0.68,0.00,0.00}{#1}}
\newcommand{\DecValTok}[1]{\textcolor[rgb]{0.68,0.00,0.00}{#1}}
\newcommand{\DocumentationTok}[1]{\textcolor[rgb]{0.37,0.37,0.37}{\textit{#1}}}
\newcommand{\ErrorTok}[1]{\textcolor[rgb]{0.68,0.00,0.00}{#1}}
\newcommand{\ExtensionTok}[1]{\textcolor[rgb]{0.00,0.23,0.31}{#1}}
\newcommand{\FloatTok}[1]{\textcolor[rgb]{0.68,0.00,0.00}{#1}}
\newcommand{\FunctionTok}[1]{\textcolor[rgb]{0.28,0.35,0.67}{#1}}
\newcommand{\ImportTok}[1]{\textcolor[rgb]{0.00,0.46,0.62}{#1}}
\newcommand{\InformationTok}[1]{\textcolor[rgb]{0.37,0.37,0.37}{#1}}
\newcommand{\KeywordTok}[1]{\textcolor[rgb]{0.00,0.23,0.31}{\textbf{#1}}}
\newcommand{\NormalTok}[1]{\textcolor[rgb]{0.00,0.23,0.31}{#1}}
\newcommand{\OperatorTok}[1]{\textcolor[rgb]{0.37,0.37,0.37}{#1}}
\newcommand{\OtherTok}[1]{\textcolor[rgb]{0.00,0.23,0.31}{#1}}
\newcommand{\PreprocessorTok}[1]{\textcolor[rgb]{0.68,0.00,0.00}{#1}}
\newcommand{\RegionMarkerTok}[1]{\textcolor[rgb]{0.00,0.23,0.31}{#1}}
\newcommand{\SpecialCharTok}[1]{\textcolor[rgb]{0.37,0.37,0.37}{#1}}
\newcommand{\SpecialStringTok}[1]{\textcolor[rgb]{0.13,0.47,0.30}{#1}}
\newcommand{\StringTok}[1]{\textcolor[rgb]{0.13,0.47,0.30}{#1}}
\newcommand{\VariableTok}[1]{\textcolor[rgb]{0.07,0.07,0.07}{#1}}
\newcommand{\VerbatimStringTok}[1]{\textcolor[rgb]{0.13,0.47,0.30}{#1}}
\newcommand{\WarningTok}[1]{\textcolor[rgb]{0.37,0.37,0.37}{\textit{#1}}}

\usepackage{longtable,booktabs,array}
\newcounter{none} % for unnumbered tables
\usepackage{calc} % for calculating minipage widths
% Correct order of tables after \paragraph or \subparagraph
\usepackage{etoolbox}
\makeatletter
\patchcmd\longtable{\par}{\if@noskipsec\mbox{}\fi\par}{}{}
\makeatother
% Allow footnotes in longtable head/foot
\IfFileExists{footnotehyper.sty}{\usepackage{footnotehyper}}{\usepackage{footnote}}
\makesavenoteenv{longtable}
\usepackage{graphicx}
\makeatletter
\newsavebox\pandoc@box
\newcommand*\pandocbounded[1]{% scales image to fit in text height/width
  \sbox\pandoc@box{#1}%
  \Gscale@div\@tempa{\textheight}{\dimexpr\ht\pandoc@box+\dp\pandoc@box\relax}%
  \Gscale@div\@tempb{\linewidth}{\wd\pandoc@box}%
  \ifdim\@tempb\p@<\@tempa\p@\let\@tempa\@tempb\fi% select the smaller of both
  \ifdim\@tempa\p@<\p@\scalebox{\@tempa}{\usebox\pandoc@box}%
  \else\usebox{\pandoc@box}%
  \fi%
}
% Set default figure placement to htbp
\def\fps@figure{htbp}
\makeatother


% definitions for citeproc citations
\NewDocumentCommand\citeproctext{}{}
\NewDocumentCommand\citeproc{mm}{%
  \begingroup\def\citeproctext{#2}\cite{#1}\endgroup}
\makeatletter
 % allow citations to break across lines
 \let\@cite@ofmt\@firstofone
 % avoid brackets around text for \cite:
 \def\@biblabel#1{}
 \def\@cite#1#2{{#1\if@tempswa , #2\fi}}
\makeatother
\newlength{\cslhangindent}
\setlength{\cslhangindent}{1.5em}
\newlength{\csllabelwidth}
\setlength{\csllabelwidth}{3em}
\newenvironment{CSLReferences}[2] % #1 hanging-indent, #2 entry-spacing
 {\begin{list}{}{%
  \setlength{\itemindent}{0pt}
  \setlength{\leftmargin}{0pt}
  \setlength{\parsep}{0pt}
  % turn on hanging indent if param 1 is 1
  \ifodd #1
   \setlength{\leftmargin}{\cslhangindent}
   \setlength{\itemindent}{-1\cslhangindent}
  \fi
  % set entry spacing
  \setlength{\itemsep}{#2\baselineskip}}}
 {\end{list}}
\usepackage{calc}
\newcommand{\CSLBlock}[1]{\hfill\break\parbox[t]{\linewidth}{\strut\ignorespaces#1\strut}}
\newcommand{\CSLLeftMargin}[1]{\parbox[t]{\csllabelwidth}{\strut#1\strut}}
\newcommand{\CSLRightInline}[1]{\parbox[t]{\linewidth - \csllabelwidth}{\strut#1\strut}}
\newcommand{\CSLIndent}[1]{\hspace{\cslhangindent}#1}



\setlength{\emergencystretch}{3em} % prevent overfull lines

\providecommand{\tightlist}{%
  \setlength{\itemsep}{0pt}\setlength{\parskip}{0pt}}



 


\KOMAoption{captions}{tableheading}
\makeatletter
\@ifpackageloaded{tcolorbox}{}{\usepackage[skins,breakable]{tcolorbox}}
\@ifpackageloaded{fontawesome5}{}{\usepackage{fontawesome5}}
\definecolor{quarto-callout-color}{HTML}{909090}
\definecolor{quarto-callout-note-color}{HTML}{0758E5}
\definecolor{quarto-callout-important-color}{HTML}{CC1914}
\definecolor{quarto-callout-warning-color}{HTML}{EB9113}
\definecolor{quarto-callout-tip-color}{HTML}{00A047}
\definecolor{quarto-callout-caution-color}{HTML}{FC5300}
\definecolor{quarto-callout-color-frame}{HTML}{acacac}
\definecolor{quarto-callout-note-color-frame}{HTML}{4582ec}
\definecolor{quarto-callout-important-color-frame}{HTML}{d9534f}
\definecolor{quarto-callout-warning-color-frame}{HTML}{f0ad4e}
\definecolor{quarto-callout-tip-color-frame}{HTML}{02b875}
\definecolor{quarto-callout-caution-color-frame}{HTML}{fd7e14}
\makeatother
\makeatletter
\@ifpackageloaded{caption}{}{\usepackage{caption}}
\AtBeginDocument{%
\ifdefined\contentsname
  \renewcommand*\contentsname{Table of contents}
\else
  \newcommand\contentsname{Table of contents}
\fi
\ifdefined\listfigurename
  \renewcommand*\listfigurename{List of Figures}
\else
  \newcommand\listfigurename{List of Figures}
\fi
\ifdefined\listtablename
  \renewcommand*\listtablename{List of Tables}
\else
  \newcommand\listtablename{List of Tables}
\fi
\ifdefined\figurename
  \renewcommand*\figurename{Figure}
\else
  \newcommand\figurename{Figure}
\fi
\ifdefined\tablename
  \renewcommand*\tablename{Table}
\else
  \newcommand\tablename{Table}
\fi
}
\@ifpackageloaded{float}{}{\usepackage{float}}
\floatstyle{ruled}
\@ifundefined{c@chapter}{\newfloat{codelisting}{h}{lop}}{\newfloat{codelisting}{h}{lop}[chapter]}
\floatname{codelisting}{Listing}
\newcommand*\listoflistings{\listof{codelisting}{List of Listings}}
\makeatother
\makeatletter
\makeatother
\makeatletter
\@ifpackageloaded{caption}{}{\usepackage{caption}}
\@ifpackageloaded{subcaption}{}{\usepackage{subcaption}}
\makeatother
\makeatletter
\@ifpackageloaded{sidenotes}{}{\usepackage{sidenotes}}
\@ifpackageloaded{marginnote}{}{\usepackage{marginnote}}
\makeatother
\usepackage{bookmark}
\IfFileExists{xurl.sty}{\usepackage{xurl}}{} % add URL line breaks if available
\urlstyle{same}
\hypersetup{
  pdftitle={State-of-the-Art Orbit Uncertainty Propagation Methods used in Space Situational Awareness (SDA)},
  colorlinks=true,
  linkcolor={blue},
  filecolor={Maroon},
  citecolor={Blue},
  urlcolor={Blue},
  pdfcreator={LaTeX via pandoc}}


\title{State-of-the-Art Orbit Uncertainty Propagation Methods used in
Space Situational Awareness (SDA)}
\author{Srihari S}
\date{}
\begin{document}
\maketitle


\section{Introduction}\label{introduction}

In the contemporary landscape of Space Domain Awareness (SDA) and Space
Situational Awareness (SSA), the precise characterization of orbital
state uncertainty is not merely a mathematical exercise but a
requirement for the preservation of the orbital environment. As the
population of resident space objects (RSOs) grows exponentially, the
risk of catastrophic collisions---such as the 2009 Iridium-Cosmos
event---escalates, potentially precipitating the ``Kessler Syndrome'',
according to which, as the density of objects in orbit increases, the
likelihood of collisions grows, leading to a cascading chain reaction
where each collision generates new debris. This self-perpetuating
process could eventually render certain orbits unusable for satellites
and spacecraft, posing a severe threat to both current and future space
operations.

Uncertainty Propagation (UP) serves as the mathematical engine for
several mission-critical tasks:

\begin{itemize}
\item
  \textbf{Conjunction Assessment (CA):} Predicting the probability of
  collision (\(P_C\)) between two objects by propagating their state
  uncertainties to the time of closest approach (TCA).
\item
  \textbf{Object Custody:} Establishing and maintaining the identity and
  tracking of newly detected objects through Bayesian filtering.
\item
  \textbf{Sensor Tasking:} Optimizing the use of limited ground-based
  radars and optical sensors by predicting where an object's probability
  mass will be at future epochs.
\item
  \textbf{Maneuver Planning:} Assessing the necessity and efficacy of
  collision avoidance maneuvers (CAMs) under non-Gaussian uncertainty
  regimes.
\end{itemize}

\subsection{Problem Formulation of Uncertainty
Propagation}\label{problem-formulation-of-uncertainty-propagation}

The stochastic uncertainty propagation problem is defined on a
probability space \((\mathcal{S}, \Sigma, \mathcal{P})\) where
\(\mathcal{S}\) is the sample space, \(\Sigma\) is the σ-algebra, and
\(\mathcal{P}\) is the probability measure.

\marginnote{\begin{footnotesize}

\begin{itemize}
\item
  \textbf{Sample Space (\(\mathcal{S}\))}: The set of all possible
  outcomes of a random experiment.\\
  \emph{Example}: For a coin toss,
  \(\mathcal{S} = \{\text{Heads}, \text{Tails}\}\).
\item
  \textbf{σ-Algebra (\(\Sigma\) )}: A collection of subsets of the
  sample space that includes \(\mathcal{S}\) and is closed under
  complements and countable unions. These subsets are the ``events'' we
  can assign probabilities to.\\
  \emph{Example}: For a die roll, \(\Sigma\) might include sets like
  \(\{1,2,3\}\), \(\{4,5,6\}\), \(\mathcal{S}\), and \(\emptyset\).
\item
  \textbf{Probability Measure (\(\mathcal{P}\))}: A function that
  assigns a number between 0 and 1 to each event in the σ-algebra,
  satisfying:

  \begin{itemize}
  \item
    \(\mathcal{P}(\mathcal{S}) = 1\)
  \item
    If events are disjoint,
    \(\mathcal{P}(\bigcup_i A_i) = \sum_i \mathcal{P}(A_i)\).
    \emph{Example}: For a fair die,
    \(\mathcal{P}(\{2\}) = \tfrac{1}{6}\).
  \end{itemize}
\end{itemize}

\end{footnotesize}}

The random input vector \(\mathbf{\xi} \in \mathbb{R}^d\) is defined on
\((\mathcal{S}, \Sigma, \mathcal{P})\) and consists of \(d\) independent
random inputs that characterize the stochastic system. These inputs are
drawn from a probability distribution defined over the space
\(\Gamma^d \subset \mathbb{R}^d\), which is the support of the random
variables. The support space \(\Gamma^d\) depends on the problem; for
example, \(\Gamma^d = [0,1]^d\) for \(d\)-dimensional uniform random
variables in the range \([0,1]\).

The uncertainty propagation problem seeks to generate a solution to a
stochastic ordinary differential equation:

\begin{equation}\protect\phantomsection\label{eq-stochastic-ode}{
\mathcal{A}(t, \mathbf{\xi}; \mathbf{x}) = 0, 
\quad (t, \mathbf{\xi}) \in [0, t_f] \times \Gamma^d, 
\quad \mathcal{P}\text{-a.s. in } \mathcal{S}
}\end{equation}

where \(t \in [0, t_f]\) is the time variable,
\(\mathbf{x} \in \mathbb{R}^n\) is a vector of Quantities of Interest
(QoIs), and \(\mathcal{A}\) is the stochastic ODE operator that defines
the flow of \(\mathbf{x}\) to \(t\) as a function of \(\mathbf{\xi}\). A
solution to the uncertainty propagation problem quantifies the effects
of random inputs \(\mathbf{\xi}\) on the random QoIs \(\mathbf{x}\) as a
function of time \(t\) and the random inputs \(\boldsymbol{\xi}\).

The random inputs \(\boldsymbol{\xi}\) may include initial state
uncertainties, force-modeling errors, and environmental perturbations.

\subsection{Taxonomy of Orbital
Uncertainties}\label{taxonomy-of-orbital-uncertainties}

The literature categorize the non-deterministic inputs into two distinct
mathematical and physical frameworks:

\begin{itemize}
\item
  \textbf{Aleatory Uncertainty (Random):}

  \begin{itemize}
  \item
    \textbf{Definition:} Characterizes the inherent, irreducible
    randomness within the physical system.
  \item
    \textbf{Examples:} Sensor measurement noise from ground-based radars
    or optical telescopes, and fluctuations in atmospheric density that
    affect drag in a stochastic manner.
  \item
    \textbf{Modeling:} Almost universally represented using precise
    probability distributions (PDFs).
  \end{itemize}
\item
  \textbf{Epistemic Uncertainty (Systematic):}

  \begin{itemize}
  \item
    \textbf{Definition:} Arises from a lack of knowledge or limited
    information about the system's physics or model parameters.
  \item
    \textbf{Examples:} Modeling inadequacies such as truncated gravity
    field harmonics (e.g., using a 2x2 vs.~70x70 model), unknown
    area-to-mass ratios for solar radiation pressure, and errors
    introduced by using fast surrogate models (metamodels) instead of
    high-fidelity propagators.
  \item
    \textbf{Modeling:} While often forced into a PDF-based
    representation, recent research suggests using Outer Probability
    Measures (OPMs) and credibilistic filtering to bound these errors
    without over-characterizing them (Jones et al. 2019).
  \end{itemize}
\end{itemize}

\subsection{Foundational Assumptions in Orbital Uncertainty
Propagation}\label{foundational-assumptions-in-orbital-uncertainty-propagation}

To make the complex dynamics of orbital motion computationally
tractable, researchers generally operate under a standardized set of
mathematical and physical assumptions:

\begin{itemize}
\item
  \textbf{Point Mass Dynamics:} For most scenarios, space objects are
  approximated as point masses, although this assumption may be relaxed
  for large structures like the International Space Station (ISS).
\item
  \textbf{State Representation:} In orbital mechanics, the QoIs are
  represented by the \textbf{state vector} \(\mathbf{x}\), representing
  position and velocity in an inertial frame (e.g., ECI) or a
  specialized coordinate set like the Modified Equinoctial Orbital
  Elements (MEqOE).
\item
  \textbf{Initial State Quantification:} It is assumed that an initial
  state \(\mathbf{x}_0\) is known at epoch \(t_0\), accompanied by a
  multivariate Probability Density Function (PDF), often initially
  modeled as Gaussian \(N(\boldsymbol{\mu}_0, \mathbf{P}_0)\), where
  \(\boldsymbol{\mu}_0=\mathbf{x}_0\) is the mean state vector and
  \(\mathbf{P}_0\) is the initial covariance matrix.
\item
  \textbf{Deterministic Flow:} The trajectory follows non-linear
  Ordinary Differential Equations (ODEs) influenced by both central-body
  gravity and various perturbative accelerations. Hence we do not
  consider epistemic uncertainty in the dynamics model itself, but
  rather focus on the propagation of the initial state uncertainty
  through these deterministic equations.
\end{itemize}

\subsection{The Fokker-Planck
Equation}\label{the-fokker-planck-equation}

The \textbf{Fokker-Planck Equation} (FPE) is the fundamental partial
differential equation (PDE) governing the time evolution of the PDF of a
dynamical system subject to both deterministic forces and stochastic
process noise.

\begin{itemize}
\item
  \textbf{Role in UP:} It describes how a probability cloud ``flows''
  and ``spreads'' through state space. The equation consists of a
  \textbf{drift term} (representing deterministic motion like Keplerian
  flow) and a \textbf{diffusion term} (representing the broadening of
  the PDF due to stochastic forcing).
\item
  \textbf{Mathematical Expression:} For a state vector \(\mathbf{x}\),
  the FPE is given by:
  \begin{equation}\protect\phantomsection\label{eq-fokker-planck}{
  \frac{\partial p(\mathbf{x}, t)}{\partial t} = -\nabla \cdot (\mathbf{f} p) + \frac{1}{2} \text{Tr}(\mathbf{\Sigma} \nabla^2 p)
  }\end{equation}

  where \(\mathbf{f}\) is the deterministic flow vector,
  \(\mathbf{\Sigma}\) is the diffusion tensor, and \(p(\mathbf{x}, t)\)
  is the PDF at time \(t\).
\item
  \textbf{The Computational Challenge in astrodynamics:} Direct
  numerical solutions to the FPE are generally intractable in
  astrodynamics due to the \textbf{``curse of dimensionality''}; a 6D
  orbital state requires spatial grids that grow exponentially,
  necessitating the use of the state-of-the-art approximation methods
  discussed in this report.
\end{itemize}

The following sections provide a detailed overview of these methods,
highlighting their mathematical and theoretical foundations, practical
implementations and variations among them, an overview on the
combinations with other major methods and performance characteristics.

\section{The Sampling Baseline: Monte Carlo
(MC)}\label{the-sampling-baseline-monte-carlo-mc}

The Monte~Carlo method serves as the ``ground truth'' and absolute
benchmark for all other techniques. The process involves generating a
large set of random samples (\(N_{\text{samples}}\)) from an initial
multivariate Gaussian distribution using the square‑right matrix (\(S\))
of the covariance matrix.

\begin{equation}\protect\phantomsection\label{eq-covariance-matrix}{
P = SS^{T}
}\end{equation}

\marginnote{\begin{footnotesize}

The covariance matrix \(P\) captures the uncertainty in the state
estimate which we get from orbit determination. The square‑right matrix
\(S\) is obtained through Cholesky decomposition or Singular Value
Decomposition (SVD) of \(P\). This allows us to generate samples that
are consistent with the initial uncertainty.

\end{footnotesize}}

The samples are created as follows:
\begin{equation}\protect\phantomsection\label{eq-mc-sampling}{
\mathbf{x}_0(\boldsymbol{\xi}) = \boldsymbol{\mu}_0 + S\boldsymbol{\xi}, \quad \boldsymbol{\xi} \sim N(0, I_6)
}\end{equation}

Where \(\boldsymbol{\xi}\) is a vector of standard normal random
variable inputs, generated using a pseudo random number simulator. The
random inputs may include initial state uncertainties, force-modeling
errors, and environmental perturbations.

\subsection{Mechanics of the Sampling-Based
Approach}\label{mechanics-of-the-sampling-based-approach}

MC simulation functions by drawing a large ensemble of random samples
from an initial distribution, which can be Gaussian or any other known
PDF.

\begin{itemize}
\item
  \textbf{Independent Propagation:} Each sample represents a potential
  realization of the system state and is propagated independently
  through numerical integration of the dynamics.
\item
  \textbf{Non-Intrusive Nature:} Because each sample is treated
  separately, MC is intrinsically non-intrusive, meaning it treats the
  dynamics model as a ``black box'' and requires no modifications to
  existing dynamics code.
\item
  \textbf{Statistical Reconstruction:} After propagation, the ensemble
  of final states provides a discrete representation of the final PDF.
  Statistical properties, such as the mean and covariance, are estimated
  directly by computing the average and sample-covariance of the
  propagated points.
\end{itemize}

\subsection{Role as the ``Ground Truth''
Baseline}\label{role-as-the-ground-truth-baseline}

The sources {[}Bos (2025){]}(Soman 2025) identify MC simulation as the
\textbf{definitive benchmark} for evaluating the accuracy of all other
UP methods.

\begin{itemize}
\tightlist
\item
  \textbf{Reliability:} MC results approach the ``true'' probability
  distribution as the number of samples goes to infinity.
\item
  \textbf{Accuracy Metric:} In comparative studies, researchers
  typically use a large-scale MC run (such as \(10^5\) samples) as the
  reference truth to calculate accuracy metrics like the \textbf{N L2
  distance} for alternative methods.
\end{itemize}

\subsection{The Computational
Trade-off}\label{the-computational-trade-off}

While highly accurate, the primary limitation of MC is its
\textbf{extreme computational intensity}.

\begin{itemize}
\tightlist
\item
  \textbf{Convergence Rate:} The rate of statistical convergence is
  proportional to the inverse square root of the ensemble size
  (\(\text{Monte Carlo Standard error}\propto 1/\sqrt{N}\), where \(N\)
  is the number of samples), meaning achieving high accuracy requires a
  massive number of samples.
\end{itemize}

\begin{tcolorbox}[enhanced jigsaw, arc=.35mm, bottomrule=.15mm, bottomtitle=1mm, breakable, colback=white, colbacktitle=quarto-callout-tip-color!10!white, colframe=quarto-callout-tip-color-frame, coltitle=black, left=2mm, leftrule=.75mm, opacityback=0, opacitybacktitle=0.6, rightrule=.15mm, title=\textcolor{quarto-callout-tip-color}{\faLightbulb}\hspace{0.5em}{Monte-Carlo Standard Error}, titlerule=0mm, toprule=.15mm, toptitle=1mm]

The statistical convergence rate describes how quickly the Monte Carlo
estimate approaches the true value as the number of samples increases.

For this, we assume that the samples are independent and identically
distributed (i.i.d.) random variables drawn from the underlying
probability distribution. Let these samples be denoted as
\(X_1, X_2, \ldots, X_N\), where \(N\) is the total number of samples.

Now we assume \(Y=f(\mathbf{X})\) to be a random variable with

\[
E[Y]=\mu,\qquad \operatorname{Var}(Y)=\sigma^2.
\]

Which gives us \(Y_1, Y_2, \ldots, Y_N\) as i.i.d. random variables
which form the output underlying uncertainty distribution. The Monte
Carlo estimator for the expected value of \(Y\) is given by,

\[
\hat{\mu}=\frac{1}{N}\sum_{i=1}^{N}Y_i.
\]

Since the samples are independent,

\[
\begin{aligned}
\operatorname{Var}(\hat{\mu})&= \operatorname{Var}\left(\frac{1}{N}\sum_{i=1}^{N}Y_i\right)\\&=\frac{1}{N^2}\sum_{i=1}^{N}\operatorname{Var}(Y_i)\frac{\sigma^2}{N}
\end{aligned}
\]

Hence, the standard error of the Monte Carlo estimator is

\[
\boxed{
\operatorname{SE}(\hat{\mu})=\sqrt{\operatorname{Var}(\hat{\mu})}\frac{\sigma}{\sqrt{N}}.
}
\]

Therefore, the statistical error of Monte Carlo decreases as

\[
\boxed{
\operatorname{SE}(\hat{\mu})=O(N^{-1/2}),
}
\]

implying that reducing the error by a factor of (k) requires
approximately (k\^{}2) times more samples.

\end{tcolorbox}

\begin{itemize}
\item
  \textbf{Collision Assessment Barriers:} For rare-event quantification,
  such as conjunction assessment where collision probabilities (\(P_C\))
  may be \(10^{-4}\) or lower, the number of required samples can reach
  \textbf{\(10^7\) or more} to achieve low statistical error.
\item
  \textbf{Practical Infeasibility:} Due to these requirements, the
  sources note that standard MC simulations are often infeasible for
  practical, real-time application to the vast population of resident
  space objects (RSOs).
\end{itemize}

\subsection{MC in the Landscape of Sampling
Methods}\label{mc-in-the-landscape-of-sampling-methods}

Current literature place MC within a broader hierarchy of sampling and
sample-reliant methods designed to mitigate its costs:

\begin{itemize}
\item
  \textbf{Deterministic Sampling (Unscented Transform):} Unlike MC's
  random sampling, the \textbf{Unscented Transform (UT)} uses a small,
  deterministically chosen set of weighted \textbf{sigma points} to
  represent the first two moments of a distribution.
\item
  \textbf{Multi-Fidelity (MF) Architectures:} This ``promising''
  alternative leverages sampling concepts by propagating a large set of
  samples through fast, \textbf{low-fidelity dynamics} and then using a
  tiny subset of ``important samples'' propagated through high-fidelity
  dynamics to correct the results. MF can achieve speed-ups of up to
  \textbf{4 orders of magnitude} over standard MC while maintaining a
  high level of agreement.
\item
  \textbf{Enhanced Sampling Techniques:} To improve efficiency,
  researchers utilize variance reduction techniques like
  \textbf{importance sampling} or \textbf{stratified sampling}, as well
  as \textbf{Quasi-Monte Carlo} methods which use special sequences for
  more uniform coverage of the uncertainty space.
\item
  \textbf{Parallelization:} Leveraging multicore CPUs or GPU
  acceleration is highlighted as a primary way to reduce the run times
  of large-scale Monte Carlo campaigns.
\end{itemize}

\section{Linearized Covariance (LinCov)}\label{sec-lincov}

In the taxonomy of uncertainty propagation (UP) methods,
\textbf{Linearised Covariance (LinCov)} represents the foundational
entry in the category of \textbf{Linear Approximation Methods}. It is
characterized by its reliance on first-order Taylor-series expansions to
map statistical moments through nonlinear dynamical systems. While
computationally superior in terms of execution time, its mathematical
architecture introduces significant epistemic risks when the system's
nonlinearity induces non-Gaussianity.

\subsection{Mathematical Framework and
Mechanics}\label{mathematical-framework-and-mechanics}

LinCov operates on the fundamental assumption that if a reference
trajectory \(x^*(t)\) is sufficiently close to the true trajectory
\(x(t)\), the deviation \(\Delta x(t)\) can be modeled as a linear
system.

\begin{itemize}
\item
  \textbf{First-Order Truncation:} The method defines the system
  dynamics \(\dot{x}(t) = F(x, t)\) and expands them into a Taylor
  series about the reference trajectory. By neglecting higher-order
  terms (second-order and above), the dynamics are simplified to
  \(\Delta\dot{x}(t) \approx J^*(t)\Delta x(t)\), where \(J^*\) is the
  Jacobian matrix evaluated along the reference path.
\item
  \textbf{The State Transition Matrix (STM):} The solution to this
  linear system is facilitated by the STM, \(\Phi(t, t_0)\), which maps
  the initial deviation to a future epoch.
\end{itemize}

\begin{tcolorbox}[enhanced jigsaw, arc=.35mm, bottomrule=.15mm, bottomtitle=1mm, breakable, colback=white, colbacktitle=quarto-callout-tip-color!10!white, colframe=quarto-callout-tip-color-frame, coltitle=black, left=2mm, leftrule=.75mm, opacityback=0, opacitybacktitle=0.6, rightrule=.15mm, title=\textcolor{quarto-callout-tip-color}{\faLightbulb}\hspace{0.5em}{State Transition Matrix (STM)}, titlerule=0mm, toprule=.15mm, toptitle=1mm]

The STM is formally defined as the solution to the matrix differential
equation \(\dot{\Phi}(t, t_0) = \mathbf{J}^*(t)\Phi(t, t_0)\), with the
initial condition \(\Phi(t_0, t_0) = \mathbf{I}\) (the identity matrix).

\end{tcolorbox}

\begin{itemize}
\item
  \textbf{Moment Propagation:} The mean state vector and covariance
  matrix are defined as the expected value of the state, and the
  expected squared deviation from the mean state, respectively: \[
  \begin{aligned}
  \bar{\mathbf{x}} &= \int_{\infty} \xi p(\xi)\,d\xi \\
  \mathbf{P} &= \int_{\infty} (\xi - \bar{x})(\xi - \bar{x})^T p(\xi,t)\,d\xi
  \end{aligned}
  \]

  Assuming that the differences between the true and reference
  trajectories are small, the mean and covariance can be propagated
  through the STM as follows:
  \begin{equation}\protect\phantomsection\label{eq-lincov-propagation}{
  \begin{aligned}
  \bar{\mathbf{x}}(t) &= \Phi(t, t_0)\bar{\mathbf{x}}(t_0) \\
  \mathbf{P}(t) &= \Phi(t, t_0)\mathbf{P}(t_0)\Phi^T(t, t_0) + G(t)Q(t)G^T(t)
  \end{aligned}
  }\end{equation}

  Where \(G(t)\) is the process noise input matrix and \(Q(t)\) is the
  process noise covariance, which are used to account for force-model
  mismodelling perturbation errors. Bos (2025) does not include process
  noise in his LinCov implementation.
\item
  \textbf{Gaussian Assumption:} A defining characteristic of LinCov is
  the preservation of the distribution's functional form; it assumes
  that an initially Gaussian PDF remains Gaussian throughout the
  propagation horizon.
\end{itemize}

\subsubsection{Algorithmic Implementation of
LinCov}\label{algorithmic-implementation-of-lincov}

The following flowchart delineates the procedural logic for propagating
orbital uncertainty using the LinCov framework:

\begin{Shaded}
\begin{Highlighting}[]
\NormalTok{graph TD}
\NormalTok{    A[Initial Mean x0 and Covariance P0] {-}{-}\textgreater{} B[Define Reference Trajectory x*]}
\NormalTok{    B {-}{-}\textgreater{} C[Integrate Nonlinear Dynamics for x*]}
\NormalTok{    C {-}{-}\textgreater{} D[Compute Time{-}Varying Jacobian J*]}
\NormalTok{    D {-}{-}\textgreater{} E[Solve Variational Equations for STM: ¤Φ = J*Φ]}
\NormalTok{    E {-}{-}\textgreater{} F[Map Final Mean: x\_tf = Φ x0]}
\NormalTok{    F {-}{-}\textgreater{} G[Map Final Covariance: P\_tf = Φ P0 ΦT]}
\NormalTok{    G {-}{-}\textgreater{} H[Output Analytic Gaussian PDF]}
\end{Highlighting}
\end{Shaded}

\subsection{Performance Analysis: The Efficiency-Accuracy
Trade-off}\label{performance-analysis-the-efficiency-accuracy-trade-off}

LinCov's role in Space Domain Awareness (SDA) is defined by its extreme
computational efficiency, often being several orders of magnitude faster
than sampling-based methods.

\begin{itemize}
\item
  \textbf{Computational Primacy:} LinCov is the fastest implemented
  method in modern comparative studies because it requires the
  integration of only one nominal state vector to solve the variational
  equations.
\item
  \textbf{Failure in Nonlinear Regimes:} The accuracy of LinCov
  inversely correlates with propagation time and dynamical nonlinearity
  (e.g., highly elliptical orbits). As the true PDF warps into
  ``banana'' or ``kidney'' shapes, the linear mapping fails to capture
  the curvature, leading to a significant increase in the \(NL_2\)
  distance metric relative to the Monte Carlo truth.
\item
  \textbf{Structural Limitations:} While LinCov captures the location
  (mean) and the general spread (scale) of the uncertainty, it
  systematically ignores the higher-order moments (skewness and
  kurtosis) that define the ``tails'' of a non-Gaussian distribution.
\end{itemize}

\subsection{Contextual Nuance: Surprising Utility in Conjunction
Assessment}\label{contextual-nuance-surprising-utility-in-conjunction-assessment}

A critical insight provided by Bos (2025) is that LinCov's failure to
approximate the \emph{shape} of a PDF does not necessarily preclude its
utility in specific \textbf{Conjunction Assessment (CA)} tasks.

\begin{itemize}
\item
  \textbf{Preservation of Decision Metrics:} In simulations involving
  Low Earth Orbit (LEO) conjunctions with long propagation times (e.g.,
  96 hours), LinCov produced collision probabilities (\(P_C\)) nearly
  identical to baseline Monte Carlo runs, despite the underlying
  distributions being highly non-Gaussian.
\item
  \textbf{Geometric Explanation:} This ``surprising'' performance
  suggests that for certain conjunction geometries, capturing the
  ``length'' and ``width'' of the uncertainty spread is more critical
  for \(P_C\) calculation than accurately modeling the time of closest
  approach (TCA) uncertainty distribution.
\item
  \textbf{Caveats:} This robust performance may be an artifact of
  specific scenario settings, such as identical orbital regimes for both
  objects or symmetrical initial uncertainties.
\end{itemize}

\subsection{Taxonomy and Comparison}\label{taxonomy-and-comparison}

LinCov sits alongside several related techniques that utilize varying
forms of linearization:

\begin{itemize}
\item
  \textbf{Statistical Linearization (CADET):} Unlike LinCov, the
  \textbf{Covariance Analysis Describing function Technique (CADET)}
  approximates system functions in an ``average'' way to minimize mean
  square error and does not strictly rely on the differentiability of
  the dynamics.
\item
  \textbf{State Transition Tensors (STT):} This framework extends LinCov
  into higher-order statistics; a first-order STT solution corresponds
  exactly to the STM used in standard LinCov, while higher-order tensors
  allow for the propagation of uncertainties through stronger
  nonlinearities without random sampling. See Section~\ref{sec-stt} for
  a detailed discussion.
\item
  \textbf{Coordinate Frame Synergy:} The accuracy of linear methods is
  significantly enhanced when operating in \textbf{Generalized
  Equinoctial Orbital Elements (GEqOE)} or \textbf{Modified Equinoctial
  Elements (MEE)}, which render the dynamics more nearly linear than
  Cartesian coordinates.
\end{itemize}

\section{Gaussian Mixture Models (GMM)
Framework}\label{gaussian-mixture-models-gmm-framework}

In the larger landscape of \textbf{Uncertainty Propagation (UP)},
\textbf{Gaussian Mixture Models (GMM)} serve as a versatile ``mixture
model'' framework that bridges the gap between computationally efficient
linear methods and the complex, non-Gaussian realities of orbital
mechanics. The foundational principle of the GMM framework is the
``universal approximator'' property, which asserts that any arbitrary
probability density function (PDF) can be represented with sufficient
accuracy as a weighted sum of multiple Gaussian kernels. This property
allows GMMs to capture the intricate shapes of non-Gaussian
distributions, such as the ``banana'' or ``kidney'' forms that emerge in
long-term propagation through the highly nonlinear dynamics of orbital
mechanics.

\subsection{Mathematical Architecture of the GMM
Framework}\label{mathematical-architecture-of-the-gmm-framework}

The GMM framework represents a complex, non-Gaussian state
\(\mathbf{x}\) as a linear combination of \(N_g\) multivariate Gaussian
components.

\subsubsection{The Joint GMM PDF}\label{the-joint-gmm-pdf}

The non-Gaussian PDF of a state vector \(\mathbf{x}\) is defined as: \[
p(\mathbf{x}) = \sum_{i=1}^{L} \alpha_i \mathbfcal{N}(\mathbf{x}; \hat{\mathbf{x}}_i, \mathbf{P}_i)
\]

\textbf{Terms Defined:}

\begin{itemize}
\item
  \(p(\mathbf{x})\): The total non-Gaussian probability density function
  of the state.
\item
  \(L\): The total number of Gaussian mixture elements (GMEs).
\item
  \(\alpha_i\): The scalar weight of the \(i\)-th component, where
  \(0 \le \alpha_i \le 1\) and \(\sum_{i=1}^{L} \alpha_i = 1\).
\item
  \(\mathbfcal{N}(\cdot)\): The \(i\)-th multivariate Gaussian kernel,
  defined as: \[
  \mathbfcal{N}(\mathbf{x}; \hat{\mathbf{x}}_i, \mathbf{P}_i) = \frac{1}{\sqrt{(2\pi)^n |\mathbf{P}_i|}} \exp \left( -\frac{1}{2} (\mathbf{x} - \hat{\mathbf{x}}_i)^T \mathbf{P}_i^{-1} (\mathbf{x} - \hat{\mathbf{x}}_i) \right)
  \]

  Where,

  \begin{itemize}
  \item
    \(\hat{\mathbf{x}}_i\): The mean state vector of the \(i\)-th
    Gaussian component.
  \item
    \(\mathbf{P}_i\): The covariance matrix of the \(i\)-th Gaussian
    component.
  \end{itemize}
\end{itemize}

\subsection{General Methodology for GMM
Propagation}\label{general-methodology-for-gmm-propagation}

The propagation of uncertainty via GMMs typically follows a three-stage
``split-propagate-recombine'' paradigm.

\begin{itemize}
\item
  \textbf{Initial Splitting:} An initial Gaussian distribution (often
  from Orbit Determination) is replaced by \(L\) components. The
  effectiveness of this stage is highly dependent on the choice of
  splitting directions and the number of components, which are
  influenced by the algorithms used to initialize and manage the mixture
  components:

  \begin{itemize}
  \item
    \textbf{Fixed vs.~Adaptive Splitting:}

    \begin{itemize}
    \item
      \textbf{Static Splitting:} The number of components is fixed at
      the initial epoch. This is common in the \textbf{Multidirectional
      GMM (MGMM)} approach, where a regular grid is formed by
      recursively applying the splitting library along multiple
      nonlinear directions.
    \item
      \textbf{Adaptive Splitting (AEGIS/AGM):} Splits are triggered
      online when a nonlinearity threshold is reached. The
      \textbf{Adaptive Entropy-based GMM (AEGIS)} monitors the deviation
      between linear and nonlinear differential entropy estimates to
      trigger splits.
    \end{itemize}
  \item
    \textbf{The Horwood Suboptimal Algorithm:} This specific technique
    refines a multivariate Gaussian along a single direction, which is
    computationally faster than multi-directional splitting without
    significantly compromising accuracy.
  \item
    \textbf{Variance-Preserving (Homoscedastic) Splitting:} While
    homoscedastic libraries often discard skewness, variance-preserving
    libraries are optimized under the constraint of matching the first
    two moments of the original distribution exactly. This ensures that
    the global mean and covariance of the mixture are identical to the
    original Gaussian before any nonlinear propagation occurs.
  \item
    \textbf{Multidirectional GMM (MGMM):} This approach recursively
    applies the splitting library along multiple nonlinear directions,
    forming a regular grid of components. While this can capture complex
    non-Gaussian features, it is computationally expensive and may lead
    to an exponential increase in the number of components.
  \end{itemize}
\item
  \textbf{Coordinate Selection:} Splitting in Cartesian space is often
  inefficient because nonlinearity acts in all six dimensions. Instead,
  curvilinear coordinates (e.g., Equinoctial or Delaunay elements) are
  used, as nonlinearity is primarily restricted to the mean longitude
  direction, allowing for a 1D split that avoids the ``curse of
  dimensionality''.
\item
  \textbf{Component-wise Propagation:} Each component
  \((\hat{\mathbf{x}}_i, \mathbf{P}_i)\) is evolved independently
  through the nonlinear dynamics using a chosen propagation method.
\item
  \textbf{Recombination:} At the target epoch, the weighted sum of
  components provides the final non-Gaussian PDF for tasks like
  Conjunction Assessment (CA).
\end{itemize}

\subsection{Algorithmic Methodology of Gaussian Mixture Modelling of
uncertainty}\label{algorithmic-methodology-of-gaussian-mixture-modelling-of-uncertainty}

The following flowchart illustrates the procedural logic for propagating
orbital uncertainty using the GMM framework. Here we use the
nonlinearity metric to pre-determine the splitting direction and the
number of components to generate.

\begin{Shaded}
\begin{Highlighting}[]
\NormalTok{graph TD}
\NormalTok{    A[Initial Gaussian State \& Covariance] {-}{-}\textgreater{} B\{Select Coordinate System\}}
\NormalTok{    B {-}{-} Cartesian {-}{-}\textgreater{} C[Calculate Nonlinearity Metric Stirling Interpolation]}
\NormalTok{    B {-}{-} Equinoctial/MEqOE {-}{-}\textgreater{} D[Select Dominant Direction Mean Longitude]}
\NormalTok{    C {-}{-}\textgreater{} E[Apply Splitting Library Rule 1/Homoscedastic]}
\NormalTok{    D {-}{-}\textgreater{} E}
\NormalTok{    E {-}{-}\textgreater{} F[Generate L Initial Gaussian Components]}
\NormalTok{    F {-}{-}\textgreater{} G[Propagate Each Component independently]}
\NormalTok{    G {-}{-} Using UT/STT/DA/PCE/MF {-}{-}\textgreater{} H[Sum Propagated Components via Weights]}
\NormalTok{    H {-}{-}\textgreater{} I[Final Non{-}Gaussian Posterior PDF]}
\end{Highlighting}
\end{Shaded}

\subsection{Variations and Hybrid
Architectures}\label{variations-and-hybrid-architectures}

The GMM framework is frequently combined with other UP methods to
optimize the speed-accuracy Pareto front.

\subsubsection{GMM-UT (Unscented
Transform)}\label{gmm-ut-unscented-transform}

The most common implementation, where each kernel's statistics are
propagated via \(2n+1\) sigma points.

\begin{itemize}
\item
  \textbf{Methodology:} Each component is represented by \(2n+1\)
  deterministic sigma points \(\chi_j\), which are propagated through
  the full nonlinear dynamics to reconstruct component statistics.
\item
  \textbf{Limitation:} The cost scales linearly with the number of
  components \(L\) and the state dimension \(n\), i.e.,
  \(O(L \cdot n)\).
\end{itemize}

\subsubsection{GMM-MF (Multi-Fidelity)}\label{gmm-mf-multi-fidelity}

A high-efficiency hybrid where the Multi-Fidelity method propagates the
sigma points of the GMM kernels.

\begin{itemize}
\item
  \textbf{Methodology:} The entire ensemble of sigma points from all
  components is propagated via a computationally cheap low-fidelity (LF)
  model (e.g., SGP4). A pivoted Cholesky decomposition of the LF Gramian
  identifies a subset of ``important samples'' to be corrected by the HF
  model via stochastic collocation. See Section~\ref{sec-mf-sc} for a
  detailed discussion on the stochastic collocation multi-fidelity
  method.
\item
  \textbf{Performance:} Achieves accuracy similar to GMM-UT but with an
  \textbf{80\% reduction in computation time}.
\end{itemize}

\subsubsection{GMM-STT (State Transition
Tensor)}\label{gmm-stt-state-transition-tensor}

This method regards Gaussian distributions as a function space to
express PDFs in closed form.

\begin{itemize}
\item
  \textbf{Methodology:} Instead of solving separate ODEs for each mean
  and covariance, the STT solution flow is reformulated as a matrix
  function valid over all components.
\item
  \textbf{Mathematical Formula (State Deviation):} \[
  \delta \hat{x}_{i}^k(t) \approx \sum_{p=1}^{m} \frac{1}{p!} \Phi^k_{j_1...j_p}(t, t_0) \delta \hat{x}_{0,i}^{j_1} \dots \delta \hat{x}_{0,i}^{j_p}
  \]

  \begin{itemize}
  \tightlist
  \item
    \(\delta \hat{x}_{i}^k(t)\) is the \(k\)-th state component of the
    \(i\)-th component; \(\Phi^k_{j_1...j_p}\) is the \(p\)-th order
    State Transition Tensor.
  \end{itemize}
\item
  \textbf{Efficiency:} It becomes more efficient than individual
  component propagation for \(L > 7\) components (using a 2nd-order
  expansion).
\end{itemize}

\subsubsection{GMM-DA (Differential
Algebra)}\label{gmm-da-differential-algebra}

The core innovation of GMM-DA is that the high-order Taylor expansion of
the solution flow is computed only once. The solution flow is expanded
into a polynomial up to an arbitrary order \(k\) without manual
derivation of variational equations. All individual Gaussian kernels are
then evaluated through this single polynomial map, turning the
propagation of thousands of components into a series of efficient
algebraic evaluations

\begin{itemize}
\item
  \textbf{Advantage:} Self-adaptive to any regular nonlinear system
  without requiring manual derivation of derivatives.
\item
  \textbf{Variant (DA-LOADS):} Uses a \textbf{Nonlinearity Index (NLI)}
  to trigger adaptive splitting only when the Taylor expansion's
  truncation error exceeds a threshold.
\end{itemize}

\subsubsection{GMM-PCE (Polynomial Chaos
Expansion)}\label{gmm-pce-polynomial-chaos-expansion}

This approach is analogous to \textbf{hp-refinement} in Finite Element
Methods.

\begin{itemize}
\item
  \textbf{Methodology:} Splitting (h-refinement) reduces the domain of
  approximation for each component, allowing for lower-order polynomials
  (p-refinement) to be used effectively.
\item
  \textbf{Benefit:} It reduces the factorial growth of PCE terms in
  high-dimensional states while maintaining the ability to capture full
  shapes of non-Gaussian distributions.
\end{itemize}

\subsubsection{GMM-Moment Method}\label{gmm-moment-method}

This semi-analytical framework propagates the first raw moment (mean)
and the second central moment (variance) of each Gaussian component
through a nonlinear system using the first/or second-order Taylor
expansion of the explicit flow function of the dynamical system.
Section~\ref{sec-moment-method} discusses the mathematical formulation
and implementation of the moment method in detail.

\subsection{Advanced Variations on the GMM
Method}\label{advanced-variations-on-the-gmm-method}

\subsubsection{Adaptive GMMs (AEGIS and
MDAES)}\label{adaptive-gmms-aegis-and-mdaes}

The \textbf{Adaptive Entropy-based Gaussian-mixture Information
Synthesis (AEGIS)} method does not rely on a fixed number of components.

\begin{itemize}
\item
  \textbf{Mechanism:} It monitors component nonlinearity by comparing
  the linear evolution of its differential entropy against a nonlinear
  estimate (e.g., via UKF).
\item
  \textbf{Directional Splitting:} To combat the ``curse of
  dimensionality,'' GMM splitting is often restricted to the eigenvector
  corresponding to the maximum eigenvalue (\(V_{max}\)) or directions
  maximizing a nonlinearity metric.

  \begin{itemize}
  \tightlist
  \item
    \textbf{Trigger:} If the entropy deviation exceeds a threshold, the
    component is split online to maintain linearity.
  \end{itemize}
\item
  \textbf{Coordinate Frame Synergy:} Pairing GMMs with \textbf{Modified
  Equinoctial Elements (MEqOE)} is highly efficient, as uncertainty in
  MEqOE is dominated by a single dimension (true longitude), requiring
  fewer kernels than Cartesian space.
\item
  \textbf{Variation:} \textbf{MDAES} combines initial multi-directional
  splitting with AEGIS to prevent correlation buildup in highly deformed
  distributions.
\end{itemize}

Despite its robustness, the sources highlight critical drawbacks:

\begin{itemize}
\item
  \textbf{The ``Centred Collection'' Effect:} In Cartesian frames,
  individual GMM components can end up clumping toward the center of the
  distribution at terminal time, failing to capture the thin ``tails''
  of the true non-Gaussian distribution.
\item
  \textbf{Combinatoric Growth:} The number of components in a batch GMM
  can scale factorially with the number of measurements
  (\(N_{\nu,0} \times \dots \times N_{\nu,L}\)), necessitating
  aggressive pruning and merging techniques.
\item
  \textbf{Tail Sensitivity:} While GMMs improve upon single Gaussians,
  they may still require an intractable number of kernels to accurately
  resolve the \(10^{-5}\) or \(10^{-6}\) tail probabilities required for
  actionable conjunction screening.
\end{itemize}

\subsubsection{Multidirectional GMM
(MGMM)}\label{multidirectional-gmm-mgmm}

Standard GMMs often split along a single axis. \textbf{MGMMs}
recursively apply splitting libraries along multiple directions to form
a regular grid in probability space.

\begin{itemize}
\item
  \textbf{Nonlinearity Metric:} Directions are ranked using a
  second-order divided difference (Stirling's interpolation) to identify
  where splits provide the most accuracy gain. \[
  \phi = \frac{f(\mathbf{m} + h\sigma \mathbf{a}) + f(\mathbf{m} - h\sigma \mathbf{a}) - 2f(\mathbf{m})}{2h^2}
  \]
\item
  \textbf{Term Definitions:} \(\mathbf{a}\) is the splitting unit vector
  and \(h = \sqrt{3}\) is the recommended step size.
\end{itemize}

\subsubsection{Gaussian Mixture Batch Processor
(GMBP)}\label{gaussian-mixture-batch-processor-gmbp}

Applied to \textbf{Initial Orbit Determination (IOD)}, the GMBP
translates a batch of noisy measurements into a GMM-based PDF using
Bayes' rule.

\begin{itemize}
\tightlist
\item
  \textbf{Methodology:} It disperses linearization errors across an
  ensemble of Gaussian components, allowing for accurate reconstruction
  of ``crescent-shaped'' posteriors common in cislunar tracking.
\end{itemize}

\subsubsection{Role in Conjunction Assessment
(CA)}\label{role-in-conjunction-assessment-ca}

In the field of collision analysis, GMMs enable \textbf{all-on-all
analysis}. Rather than evaluating millions of individual sample pairs as
in Monte Carlo simulation, every Gaussian component of a primary object
is evaluated against every component of a secondary object. This allows
for the use of more advanced analytical or numerical collision
probability (\(P_c\)) computation methods on each individual component
pair.

\subsection{Adaptive splitting in the Context of Gaussian Mixture Models
(GMM)}\label{adaptive-splitting-in-the-context-of-gaussian-mixture-models-gmm}

Bos's thesis (Bos 2025) defines Adaptive Entropy-based Gaussian-mixture
Information Synthesis (AEGIS) as an advanced framework that utilizes
automatic domain splitting to maintain an accurate representation of a
probability density function (PDF) as its non-Gaussianity increases
during propagation.

While standard GMMs often approximate a non-Gaussian PDF by splitting an
initial distribution into a fixed number of weighted components at the
start of propagation (\(t_0\)), \textbf{AEGIS performs this splitting
online}. The method treats the uncertainty as a sum of Gaussian
``kernels'' that can be individually propagated using efficient
techniques like the \textbf{Unscented Transform} or \textbf{Linearised
Covariance}.

The core philosophy of AEGIS is that by repeatedly splitting a Gaussian
component into smaller sub-components when nonlinear effects become too
high, the method reduces the local impact of those nonlinearities,
allowing individual components to remain near-Gaussian for a longer
duration.

\subsubsection{Mechanics of Adaptive
Splitting}\label{mechanics-of-adaptive-splitting}

AEGIS triggers a split based on a comparison between \textbf{linear and
nonlinear differential entropy}:

\begin{itemize}
\item
  \textbf{Linear Entropy:} This is calculated by integrating the time
  derivative of the differential entropy, which is a function of the
  system's \textbf{Jacobian matrix} (the basis of the Linearized
  Covariance method).
\item
  \textbf{Nonlinear Entropy:} This is calculated directly from the
  covariance matrix after propagating the component using a nonlinear
  technique, typically the \textbf{Unscented Transform}.
\item
  \textbf{The Splitting Threshold:} When the normalized difference
  between these two entropy calculations exceeds a user-defined
  threshold (\(\epsilon_{ent}\)), the propagation is paused, the
  component is split into smaller Gaussians, and the entropies are
  recalculated before propagation continues.
\end{itemize}

\subsubsection{Performance and Computational
Challenges}\label{performance-and-computational-challenges}

Despite its theoretical potential, the study highlight significant
practical drawbacks to the AEGIS method in orbital scenarios:

\begin{itemize}
\item
  \textbf{Extreme Computational Intensity:} AEGIS was found to be the
  most computationally expensive method tested. Because it requires
  propagating the State Transition Matrix (STM) and calculating
  differential entropy at every time step, it is slower than a standard
  UT even when \textbf{no splitting occurs}.
\item
  \textbf{Exponential Component Growth:} For long propagation horizons,
  the number of components can increase exponentially as splits occur at
  increasingly frequent intervals. This leads to \textbf{unreachable
  computation times} or out-of-memory errors, making the method
  impractical for scenarios spanning many orbital revolutions.
\item
  \textbf{Accuracy Decay:} While AEGIS approximates distributions well
  for short durations (e.g., 2 revolutions), its accuracy decays sharply
  in long-term Low Earth Orbit (LEO) cases. In these scenarios, it is
  generally outperformed by the \textbf{Multi-Fidelity (MF)} and
  \textbf{Polynomial Chaos Expansion (PCE)} methods.
\end{itemize}

\subsubsection{Impact of Coordinate
Frames}\label{impact-of-coordinate-frames}

The study notes that the poor performance of AEGIS was partly due to
operating in \textbf{Cartesian coordinates}, where nonlinearity affects
all dimensions. The sources suggest that combining AEGIS with
\textbf{Modified Equinoctial Elements (MEqOE)} would be more effective;
in MEqOE, the uncertainty primarily distributes along a single direction
(true longitude), which would likely require far fewer splits and
components to achieve high accuracy.

\section{Moment Method}\label{sec-moment-method}

The \textbf{Moment Method} (or Method of Moments) is a statistical
technique utilized in uncertainty propagation to estimate the evolution
of a probability density function (PDF) by focusing on its statistical
moments---primarily the mean and covariance. Within the hierarchy of
uncertainty propagation methods, it serves as a computationally
efficient bridge between low-fidelity linear approximations and
high-fidelity, but expensive, sampling-based techniques like Monte Carlo
(MC).

\subsection{Mathematical Framework and
Methodology}\label{mathematical-framework-and-methodology}

The method relies on the assumption that if an input variable (or
vector) is characterized by a Gaussian distribution, its first two
moments are necessary and sufficient to describe its state (Soman 2025).

\subsubsection{Taylor Series Expansion}\label{taylor-series-expansion}

The core mechanism involves expanding the nonlinear system function into
a Taylor series about the mean value of the inputs.

\textbf{Explicit vs.~Implicit Formulations:}

\begin{itemize}
\item
  \textbf{Explicit Moment Method:} Applied when a direct functional
  relationship \(y = f(x)\) exists between input and output.
\item
  \textbf{Implicit Moment Method:} Utilized in astrodynamics when only
  the governing differential equations \(\dot{x} = f(x, t)\) are known.
  It requires a numerical integrator (e.g., RKF78) to define the
  solution flow \(\phi\), which is then expanded.
\end{itemize}

\subsubsection{Moment Propagation}\label{moment-propagation}

In order to explain the moment propagation, consider an explicit moment
propagation problem with a nonlinear function \(y = f(x)\) where \(x\)
is a random variable with mean \(\mu_x\) and covariance \(\sigma_x^2\).
The first two moments of the output \(y\) can be approximated depending
on the order of the Taylor approximation:

\textbf{First-Order Equations:}

\begin{enumerate}
\def\labelenumi{\arabic{enumi}.}
\item
  \textbf{Mean (\(\mu_y\)):} \(\mu_y = f(\mu_x)\)
\item
  \textbf{Variance (\(\sigma_y^2\)):}
  \(\sigma_y^2 = [f'(\mu_x)]^2 \sigma_x^2\)
\end{enumerate}

While fast, moment propagation on the first-order taylor approximation
fails to capture the ``mean shift'' caused by nonlinearities and
performs poorly when the PDF becomes highly non-Gaussian.

\textbf{Second-Order Equations:}

\begin{enumerate}
\def\labelenumi{\arabic{enumi}.}
\item
  \textbf{Mean:}
  \(\mu_y = f(\mu_x) + \frac{1}{2} f''(\mu_x) \sigma_x^2\)
\item
  \textbf{Variance:}
  \(\sigma_y^2 = [f'(\mu_x)]^2\sigma_x^2 + f'(\mu_x)f''(\mu_x)\mu_3^{(x)} + \frac{1}{4}[f''(\mu_x)]^2 (\mu_4^{(x)}-\sigma_x^4)\)
\end{enumerate}

\textbf{Terms Defined:}

\begin{itemize}
\item
  \(f(\mu_x)\): Function evaluation at the mean.
\item
  \(f'(\mu_x)\): First derivative (sensitivity) at the mean.
\item
  \(f''(\mu_x)\): Second derivative (curvature) at the mean.
\item
  \(\mu_3^{(x)}\): The third central moment of the input distribution.
\item
  \(\mu_4^{(x)}\): The fourth central moment of the input distribution.
\item
  \(\sigma_x^4\): The square of the input variance.
\end{itemize}

\textbf{Implicit Formulation:}

In high-fidelity astrodynamics, an explicit formula \(y = f(x)\) is
rarely available. Instead, the system is defined by a set of nonlinear
Ordinary Differential Equations (ODEs):
\[\dot{\mathbf{x}} = f(\mathbf{x}, t)\]

The \textbf{Implicit Moment Method} uses the \textbf{solution flow}
\(\mathbf{x}(t)=\phi(t; \mathbf{x}_0)\) obtained via numerical
integration (e.g., Runge-Kutta schemes). The methodology treats the
numerical integrator as the mapping function:

\begin{enumerate}
\def\labelenumi{\arabic{enumi}.}
\item
  \textbf{Trajectory Integration:} Propagate the nominal mean
  \(\boldsymbol{\mu}_{x0}\) to find the final state
  \(\phi(\boldsymbol{\mu}_{x0})\).
\item
  \textbf{Sensitivity Extraction:} Compute the Jacobian and Hessian of
  the flow map, often utilizing \textbf{Automatic Differentiation (AD)}
  to ensure machine-precision accuracy.
\item
  \textbf{Moment Transformation:} Apply the multivariate moment
  equations using the flow derivatives to advance the mean and
  covariance to the target epoch \(t_f\).
\end{enumerate}

In a 6-dimensional orbital state, the method propagates the mean vector
\(\mu\) and the covariance matrix \(\Sigma\). The following sections
detail the multivariate extension of the moment method.

\subsubsection{\texorpdfstring{Multivariate Extension
(\(R^n\))}{Multivariate Extension (R\^{}n)}}\label{multivariate-extension-rn}

When the input is an \(n\)-dimensional random vector
\(\mathbf{x} \sim \mathcal{N}(\boldsymbol{\mu}_x, \mathbf{\Sigma}_x)\),
the propagation involves the Jacobian matrix and Hessian tensors.

\textbf{First-Order (Linearized) Moments:}
\[\boldsymbol{\mu}_y = f(\boldsymbol{\mu}_x)\]
\[\mathbf{\Sigma}_y = \mathbf{J} \mathbf{\Sigma}_x \mathbf{J}^T\]

\textbf{Second-Order Moments (Component-wise):}
\[\mu_{y,i} = f_i(\boldsymbol{\mu}_x) + \frac{1}{2} \text{tr}(\mathbf{H}_{fi}(\boldsymbol{\mu}_x) \mathbf{\Sigma}_x)\]
\[
\Sigma_y(i, p) = (\mathbf{J} \mathbf{\Sigma}_x \mathbf{J}^T)_{ip} + \frac{1}{2} \text{tr}(\mathbf{H}_{fi} \mathbf{\Sigma}_x \mathbf{H}_{fp} \mathbf{\Sigma}_x)
\]

\textbf{Terms Defined:}

\begin{itemize}
\item
  \(\mathbf{J}\): The \(n \times n\) Jacobian matrix,
  \(\mathbf{J} = \frac{\partial f}{\partial \mathbf{x}}\big|_{\boldsymbol{\mu}_x}\).
\item
  \(\mathbf{H}_{fi}\): The Hessian matrix of the \(i\)-th component of
  the function \(f\).
\item
  \(\text{tr}(\cdot)\): The trace operator, summing the diagonal
  elements of the resulting matrix product.
\end{itemize}

\subsection{Comparative Orders of
Approximation}\label{comparative-orders-of-approximation}

The accuracy of the Moment Method is strictly tied to the order of the
Taylor expansion used:

\begin{itemize}
\item
  \textbf{First-Order (M1):}

  \begin{itemize}
  \item
    Assumes the system is locally linear.
  \item
    Propagates the mean as the function of the mean:
    \(\mu_y = f(\mu_x)\).
  \item
    While fast, it fails to capture the ``mean shift'' caused by
    nonlinearities and performs poorly when the PDF becomes highly
    non-Gaussian.
  \end{itemize}
\item
  \textbf{Second-Order (M2):}

  \begin{itemize}
  \item
    Incorporates the \textbf{Hessian matrix} (\(H\)) of the system
    dynamics.
  \item
    The mean is corrected by a term involving the trace of the Hessian
    and the covariance:
    \(\mu_{y,i} \approx f_i(\mu_x) + \frac{1}{2} \text{tr}(H_i \Sigma_x)\).
  \item
    The covariance propagation includes additional terms to account for
    the ``spreading'' effect of curvature:
    \(\Sigma_y \approx J \Sigma_x J^T + \frac{1}{2} \text{tr}(H_i \Sigma_x H_p \Sigma_x)\).
  \item
    Empirical results show that M2 consistently maintains lower error
    norms (\(L_2\) and \(L_\infty\)) than M1, especially in nonlinear
    regimes like Highly Elliptical Orbits (HEO).
  \end{itemize}
\end{itemize}

\subsection{Integration with Gaussian Mixture Models
(GMM)}\label{integration-with-gaussian-mixture-models-gmm}

In practice, the implicit moment method is often combined with
\textbf{Gaussian Mixture Models (GMM)} to handle the ``loss of
Gaussianity'' that occurs in long-term orbital propagation. Each
Gaussian component is propagated independently using the moment method,
and the final non-Gaussian PDF is reconstructed by recombining the
weighted sum of these components.

\begin{enumerate}
\def\labelenumi{\arabic{enumi}.}
\item
  \textbf{Decomposition:} The initial Gaussian uncertainty is split into
  \(N\) smaller Gaussian components (kernels) using algorithms like the
  Horwood suboptimal algorithm.
\item
  \textbf{Local Propagation:} Each individual component has a
  sufficiently small covariance to remain ``locally Gaussian,'' allowing
  the Moment Method to propagate its first two moments accurately over
  longer durations.
\item
  \textbf{Reconstruction:} The final non-Gaussian PDF is reconstructed
  by recombining the weighted sum of these propagated components.
\end{enumerate}

\subsection{Algorithm for Multivariate Implicit GMM-Moment
Propagation}\label{algorithm-for-multivariate-implicit-gmm-moment-propagation}

The following flowchart outlines the computational algorithm for a
multivariate system (e.g., orbital mechanics in \(R^6\)) using the
implicit second-order moment method:

\begin{Shaded}
\begin{Highlighting}[]
\NormalTok{graph TD}
\NormalTok{    A[Initial State: Mean μ₀, Covariance Σ₀] {-}{-}\textgreater{} B[Horwood Splitting Algorithm]}
\NormalTok{    B {-}{-}\textgreater{} C\{For each component i=1 to N\}}
\NormalTok{    C {-}{-}\textgreater{} D[Numerical Integration: RK8/RKF78]}
\NormalTok{    D {-}{-}\textgreater{} E[Automatic Differentiation: TAPENADE]}
\NormalTok{    E {-}{-}\textgreater{} F[Extract Jacobian J and Hessian H]}
\NormalTok{    F {-}{-}\textgreater{} G[Implicit Moment Method Update]}
\NormalTok{    G {-}{-}\textgreater{} H[Propagated Component: μ\_fi, Σ\_fi]}
\NormalTok{    H {-}{-}\textgreater{} I[Recombine Components with Weights w\_i]}
\NormalTok{    I {-}{-}\textgreater{} J[Final Reconstructed PDF]}
\NormalTok{    J {-}{-}\textgreater{} K[Compute Overall Mean μ\_f and Covariance Σ\_f]}
\end{Highlighting}
\end{Shaded}

\subsection{Critical Strengths and
Limitations}\label{critical-strengths-and-limitations}

\begin{itemize}
\item
  \textbf{Computational Efficiency:} The GMM-Moment method is
  significantly faster than Monte Carlo simulations. For a 3-day LEO
  propagation, it completed in seconds/minutes compared to the hours
  required for a \(10^7\)-sample MC run.
\item
  \textbf{Derivative Sensitivity:} Unlike the Unscented Transform (UT),
  which is sample-based, the Moment Method requires explicit derivative
  information (Jacobians and Hessians).
\item
  \textbf{Automatic Differentiation (AD):} To avoid the errors of finite
  differences, modern implementations use AD (e.g., TAPENADE) to obtain
  exact derivatives by applying the chain rule to the source code of the
  propagator.
\item
  \textbf{Differentiability Requirement:} The method assumes that the
  governing equations are continuous and differentiable, which may pose
  challenges in the presence of discrete events like shadow entries (SRP
  discontinuities).
\end{itemize}

\section{Unscented Transform (UT)}\label{unscented-transform-ut}

In the rigorous taxonomy of orbital uncertainty quantification, the
\textbf{Unscented Transform (UT)} represents a pivotal
\textbf{deterministic sampling method} that bridges the gap between
computationally cheap linear approximations and prohibitively expensive
brute-force stochastic simulations. Unlike the Monte Carlo (MC) method,
which relies on random realizations, the UT operates on the foundational
intuition that it is mathematically easier to approximate a probability
distribution than it is to approximate an arbitrary nonlinear function
or transformation.

\subsection{Theoretical Context: Deterministic vs.~Random
Sampling}\label{theoretical-context-deterministic-vs.-random-sampling}

Within the landscape of sampling-based methodologies, the UT occupies a
unique niche defined by ``moderate fidelity'' and high efficiency.

\begin{itemize}
\item
  \textbf{Sigma Point Paradigm:} The UT is predicated on the intuition
  that ``it is easier to approximate a Gaussian distribution than it is
  to approximate an arbitrary nonlinear function or transformation''.
  Unlike the EKF, which linearizes the system dynamics via a first-order
  Taylor series expansion (State Transition Matrix), the UT
  parameterizes the state distribution using a set of ``sigma points''.
  While MC requires \(10^5\) to \(10^7\) random samples to achieve
  statistical convergence, the UT utilizes a minimal, carefully selected
  set of weighted \textbf{sigma points} that precisely capture the
  initial mean and covariance of a distribution.
\item
  \textbf{Black-Box Dynamics:} Like MC, the UT is non-intrusive; it
  treats the system dynamics as a ``black box,'' passing the discrete
  sigma points through the actual nonlinear equations of motion without
  requiring the derivation of Jacobians or Hessians.
\item
  \textbf{Gaussian Constraint:} A critical limitation noted in the
  sources is that the standard UT is primarily constrained by
  \textbf{Gaussian assumptions}. It nonlinearly propagates the first and
  second moments but may fail to naturally represent the
  ``banana-shaped'' non-Gaussian warpings that emerge during long-term
  orbital propagation.
\end{itemize}

\subsection{Mathematical Architecture of the
UT}\label{mathematical-architecture-of-the-ut}

For an \(n\)-dimensional state \(\mathbf{x}\) with mean
\(\bar{\mathbf{x}}\) and covariance \(\mathbf{P}_{xx}\), a set of
\(2n+1\) sigma points \(\boldsymbol{\chi}_i\) is generated using a
matrix square root (typically Cholesky decomposition) of the scaled
covariance. The mathematical selection and weight assignment are
governed by the following formulations:

\subsubsection{Sigma Point Selection}\label{sigma-point-selection}

\[\chi_0 = \bar{\mathbf{x}}\]
\begin{equation}\protect\phantomsection\label{eq-sigma-points-i}{
\chi_i = \bar{\mathbf{x}} + \left( \sqrt{(n + \kappa) \mathbf{P}} \right)_i, \quad i = 1, \dots, n
}\end{equation}
\begin{equation}\protect\phantomsection\label{eq-sigma-points-in}{
\chi_{i+n} = \bar{\mathbf{x}} - \left( \sqrt{(n + \kappa) \mathbf{P}} \right)_i, \quad i = 1, \dots, n
}\end{equation}

\subsubsection{Weight Allocation}\label{weight-allocation}

\[W_0 = \frac{\kappa}{n + \kappa}\]
\[W_i = \frac{1}{2(n + \kappa)}, \quad i = 1, \dots, 2n\]

\textbf{Terms Defined:}

\begin{itemize}
\item
  \(\bar{\mathbf{x}}\): The mean state vector.
\item
  \(\mathbf{P}\): The state covariance matrix.
\item
  \(\left( \sqrt{\cdot} \right)_i\): The \(i\)-th column of the matrix
  square root (typically calculated via \textbf{Cholesky
  decomposition}).
\end{itemize}

\begin{tcolorbox}[enhanced jigsaw, arc=.35mm, bottomrule=.15mm, bottomtitle=1mm, breakable, colback=white, colbacktitle=quarto-callout-note-color!10!white, colframe=quarto-callout-note-color-frame, coltitle=black, left=2mm, leftrule=.75mm, opacityback=0, opacitybacktitle=0.6, rightrule=.15mm, title=\textcolor{quarto-callout-note-color}{\faInfo}\hspace{0.5em}{Cholesky Decomposition}, titlerule=0mm, toprule=.15mm, toptitle=1mm]

The Cholesky decomposition of a positive definite matrix \(\mathbf{P}\)
is a lower triangular matrix \(\mathbf{L}\) such that
\(\mathbf{P} = \mathbf{L}\mathbf{L}^T\).

\end{tcolorbox}

\begin{itemize}
\tightlist
\item
  \(\kappa\): A scaling/tuning parameter, often set such that
  \(n + \kappa = 3\) for Gaussian distributions to minimize higher-order
  errors.
\end{itemize}

\subsubsection{Statistical
Reconstruction}\label{statistical-reconstruction}

After passing each \(\chi_i\) through the nonlinear dynamics
\(\mathbf{f}(\cdot)\), \[
\mathbfcal{Y}_i = \mathbf{f}(\chi_i)
\]

the terminal mean (\(\hat{\mathbf{y}}\)) and covariance
(\(\hat{\mathbf{P}}\)) are reconstructed:

\textbf{Transformed Mean:} \[
\bar{\mathbf{y}} = \sum_{i=0}^{2n} W_i \mathbfcal{Y}_i
\]

\textbf{Transformed Covariance:} \[
\bar{\mathbf{P}} = \sum_{i=0}^{2n} W_i (\mathbfcal{Y}_i - \hat{\mathbf{y}})(\mathbfcal{Y}_i - \hat{\mathbf{y}})^T
\]

\subsection{Algorithm: The Unscented Filter (UKF)
Logic}\label{algorithm-the-unscented-filter-ukf-logic}

The following flowchart delineates the procedural execution of the UT
within a predictive filtering framework, such as the Unscented Kalman
Filter:

\begin{Shaded}
\begin{Highlighting}[]
\NormalTok{graph TD}
\NormalTok{    A[Initial Mean μ and Covariance P] {-}{-}\textgreater{} B[Compute Matrix Square Root of P]}
\NormalTok{    B {-}{-}\textgreater{} C[Generate 2n+1 Sigma Points χ\_i]}
\NormalTok{    C {-}{-}\textgreater{} D[Propagate points through Nonlinear Dynamics: χ\_new = f χ\_i]}
\NormalTok{    D {-}{-}\textgreater{} E[Reconstruct Predicted Mean: μ\_pred = Σ W\_i χ\_new]}
\NormalTok{    E {-}{-}\textgreater{} F[Reconstruct Predicted Covariance: P\_pred = Σ W\_i diff\^{}2]}
\NormalTok{    F {-}{-}\textgreater{} G\{Sensor Measurement Available?\}}
\NormalTok{    G {-}{-} Yes {-}{-}\textgreater{} H[Transform Points through Observation Model: Z\_i = h χ\_new]}
\NormalTok{    H {-}{-}\textgreater{} I[Compute Innovation Covariance and Cross{-}Correlation]}
\NormalTok{    I {-}{-}\textgreater{} J[Update μ and P via Kalman Gain]}
\NormalTok{    G {-}{-} No {-}{-}\textgreater{} K[Output Propagated Analytic Gaussian PDF]}
\end{Highlighting}
\end{Shaded}

\subsection{Variations and Hybrid
Architectures}\label{variations-and-hybrid-architectures-1}

The sources identify several advanced iterations of the UT designed to
overcome specific numerical or statistical limitations:

\begin{itemize}
\item
  \textbf{Square-Root Unscented Kalman Filter (SR-UKF)} {[}Van der Merwe
  and Wan (2001){]}:

  \begin{itemize}
  \item
    \textbf{Description}: Instead of propagating the full matrix
    \(\mathbf{P}\), the SR-UKF propagates the Cholesky factor
    \(\mathbf{S}\) (where \(\mathbf{P} = \mathbf{SS}^T\)) directly.
  \item
    \textbf{Advantage}: It guarantees positive semi-definiteness of the
    covariance and offers superior numerical stability and computational
    efficiency (\(O(L^3)\) for state estimation).
  \end{itemize}
\item
  \textbf{Higher-Order Unscented Filter (HOUF)} {[}Tenne and Singh
  (2003){]}:

  \begin{itemize}
  \item
    \textbf{Description}: Expands the sigma-set to include additional
    layers of points to match input moments beyond the 4th order (e.g.,
    up to the 12th order).
  \item
    \textbf{Context}: This is essential for highly nonlinear regimes
    where standard UT truncation errors (which start at the 4th order)
    become significant.
  \end{itemize}
\item
  \textbf{Conjugate Unscented Transform (CUT)} {[}Adurthi et al.
  (2012){]}:

  \begin{itemize}
  \item
    \textbf{Description}: A non-product cubature rule that uses
    ``Conjugate Axes'' (combinations of principal axes) to satisfy
    higher-order moment equations (up to 8th order) with fewer points
    than Gauss-Hermite rules.
  \item
    \textbf{Advantage}: It captures cross-moments (e.g.,
    \(E[X_i^2 X_j^2]\)) that standard UT cannot, while avoiding the
    ``curse of dimensionality''.
  \end{itemize}
\item
  \textbf{Gaussian Mixture Model-Unscented Transform (GMM-UT)}:

  \begin{itemize}
  \item
    \textbf{Description}: To overcome the Gaussian limitation, the UT is
    frequently used to propagate the individual kernels of a \textbf{GMM
    (GMM\_UT)}. The GMM-UT method operates on the intuition that while
    the global distribution may be highly non-Gaussian (e.g., a
    ``banana'' shape), the local regions of the probability mass remain
    nearly Gaussian for a longer duration under nonlinear flow.
  \item
    \textbf{Methodology}

    \begin{itemize}
    \item
      \textbf{Sigma Point Propagation:} For each of the \(N_g\)
      components, a set of \(2n+1\) deterministic sigma points is
      generated using the component's mean \(\mathbf{x}_i\) and
      covariance \(\mathbf{P}_i\).
    \item
      \textbf{Independent Nonlinear Mapping:} These sigma points are
      passed through the high-fidelity nonlinear dynamics
      \(\mathbf{f}(\cdot)\).
    \item
      \textbf{Moment Reconstruction:} The transformed mean and
      covariance for \emph{each} kernel are reconstructed using standard
      UT weighting rules, preserving second-order accuracy.
    \item
      \textbf{Re-Synthesis:} The final non-Gaussian PDF is the
      re-weighted sum of these nonlinearly propagated Gaussian kernels
      with their corresponding weights \(\alpha_i\).
    \end{itemize}
  \end{itemize}
\item
  \textbf{AEGIS Framework:} In the \textbf{Adaptive Entropy-based
  Gaussian-mixture Information Synthesis (AEGIS)} method, the UT is
  employed to calculate ``nonlinear entropy''. By comparing this to the
  linear entropy from LinCov, the system can detect when the dynamics
  have become too nonlinear and trigger a component split.
\item
  \textbf{Multi-Fidelity (MF) Advantage:} While the UT is an efficient
  ``middle-ground,'' the sources characterize Multi-Fidelity (MF) as
  superior for high-consequence non-Gaussian cases, as MF uses
  large-scale low-fidelity sampling corrected by high-fidelity samples
  to capture complex shapes that the UT's single-Gaussian assumption
  cannot.
\end{itemize}

\subsection{Comparative Performance}\label{sec-ut-performance}

The sources highlight the UT's performance relative to other
``State-of-the-Art'' methods:

\begin{itemize}
\item
  \textbf{UT vs.~Linearised Covariance (LinCov):} The UT is consistently
  more accurate than LinCov because it captures more of the
  distribution's tails and accounts for higher-order nonlinear effects
  that first-order Taylor expansions ignore.
\item
  \textbf{UT vs.~Monte Carlo:} The UT is significantly faster than MC,
  requiring only 13 points for a standard 6D orbital state compared to
  the millions needed for MC to reach similar statistical convergence.
\end{itemize}

\section{Polynomial Chaos Expansions
(PCE)}\label{polynomial-chaos-expansions-pce}

In the rigorous domain of orbital uncertainty propagation (UP),
\textbf{Polynomial Chaos Expansion (PCE)} represents a powerful spectral
decomposition technique that characterizes the evolution of state
uncertainties through highly nonlinear dynamical systems. Unlike
traditional linearization methods that assume Gaussianity, PCE provides
a non-parametric framework capable of capturing the increasingly
non-Gaussian ``banana-shaped'' probability density functions (PDFs)
characteristic of long-term orbital motion.

\subsection{Theoretical Classification and
Role}\label{theoretical-classification-and-role}

Within the methodological landscape, PCE is distinguished by its ability
to map input uncertainties to output responses using a ``black-box''
approach that does not require alterations to existing high-fidelity
solvers.

\begin{itemize}
\item
  \textbf{Beyond Linearization:} Unlike Linearized Covariance (LinCov),
  which relies on first-order Taylor expansions and assumes Gaussian
  distributions, PCE represents inputs and outputs using a series of
  approximations that capture \textbf{higher moments} of the probability
  density function (PDF), allowing it to represent non-Gaussian shapes.
\item
  \textbf{Surrogate Modeling Paradigm:} PCE functions as a data-driven
  mathematical emulator, replacing expensive physical orbital
  propagators with a computationally cheap polynomial surrogate.
\item
  \textbf{Spectral Accuracy:} For systems where the output depends
  smoothly on the input parameters, PCE can achieve \textbf{exponential
  convergence} rates, far exceeding the \(\mathcal{O}(1/\sqrt{N})\)
  convergence of standard Monte Carlo (MC).
\end{itemize}

\subsection{Mathematical Foundations of
PCE}\label{mathematical-foundations-of-pce}

The PCE models a square-measurable stochastic output vector
\(\mathbf{X}\) as a linear combination of multivariate polynomials that
are orthogonal with respect to the joint PDF of the input uncertainties
\(\xi\).

\marginnote{\begin{footnotesize}

Here, a random variable or a stochastic solution is called
\textbf{square-measurable} if its second moment is finite:
\[\mathbb{E}[\mathbf{X}^2] < \infty\]

This condition ensures that the PCE framework can represent the random
solution as a series of orthogonal polynomials and that statistical
moments (like mean and variance) can be computed analytically from the
expansion coefficients.

\end{footnotesize}}

\subsubsection{\texorpdfstring{\textbf{The PCE Surrogate
Formula}}{The PCE Surrogate Formula}}\label{the-pce-surrogate-formula}

The stochastic output at a future time \(t\) is approximated by the
finite series expansion: \[
\hat{\mathbf{X}}(t, \xi) \approx \sum_{j=0}^{P} c_j(t) \Psi_j(\xi)
\]

\textbf{Terms Defined:}

\begin{itemize}
\item
  \(\hat{\mathbf{X}}(t, \xi)\): The approximated Quantity of Interest
  (QoI) vector (e.g., the terminal orbital state).
\item
  \(\xi\): A vector of independent standard random variables
  representing input uncertainties.
\item
  \(\Psi_j(\xi)\): Multivariate orthonormal basis functions, typically
  formed as the product of univariate polynomials.
\item
  \(c_j(t)\): Deterministic spectral coefficients (weights) determined
  at the epoch of interest.
\item
  \(P\): The number of terms in the expansion, determined by the
  polynomial order \(p\) and input dimension \(d\).
\end{itemize}

\subsubsection{\texorpdfstring{\textbf{Expansion Order and the Curse of
Dimensionality}}{Expansion Order and the Curse of Dimensionality}}\label{expansion-order-and-the-curse-of-dimensionality}

The total number of coefficients \(P\) required for a given order \(p\)
and dimension \(d\) is calculated as: \[P+1 = \frac{(p+d)!}{p!d!}\]

This factorial growth highlights the primary limitation of classical
PCE: the \textbf{curse of dimensionality}, where high-dimensional
problems (e.g., \(d > 10\)) become computationally demanding due to the
rapid growth of expansion terms.

\subsubsection{\texorpdfstring{\textbf{Orthonormality and Moment
Matching}}{Orthonormality and Moment Matching}}\label{orthonormality-and-moment-matching}

The basis functions are chosen according to the \textbf{Wiener-Askey
scheme}, ensuring they are orthogonal relative to the input PDF,
\(\rho(\xi)\): \[
\langle \Psi_\alpha, \Psi_\beta \rangle = \int_{\Gamma} \Psi_\alpha(\xi) \Psi_\beta(\xi) \rho(\xi) d\xi = \delta_{\alpha\beta}
\]

where \(\delta_{\alpha\beta}\) is the Kronecker delta. This property
allows for the \textbf{analytical evaluation of statistical moments}
directly from the spectral coefficients without additional sampling:

\begin{itemize}
\item
  \textbf{Mean:} \(\mu = c_0\).
\item
  \textbf{Variance:} \(\sigma^2 = \sum_{j=1}^P c_j^2 \|\Psi_j\|^2\).
\end{itemize}

\begin{tcolorbox}[enhanced jigsaw, arc=.35mm, bottomrule=.15mm, bottomtitle=1mm, breakable, colback=white, colbacktitle=quarto-callout-tip-color!10!white, colframe=quarto-callout-tip-color-frame, coltitle=black, left=2mm, leftrule=.75mm, opacityback=0, opacitybacktitle=0.6, rightrule=.15mm, title=\textcolor{quarto-callout-tip-color}{\faLightbulb}\hspace{0.5em}{Polynomial Families and Corresponding Input Distributions}, titlerule=0mm, toprule=.15mm, toptitle=1mm]

Usually, the choice of polynomial family is dictated by the distribution
of the input uncertainties:

\begin{itemize}
\item
  \textbf{Hermite Polynomials:} For Gaussian inputs.
\item
  \textbf{Legendre Polynomials:} For Uniform inputs.
\end{itemize}

\end{tcolorbox}

\subsection{Procedural Logic: Non-Intrusive PCE
Workflow}\label{procedural-logic-non-intrusive-pce-workflow}

Non-intrusive PCE (NIPC) is favored in astrodynamics because it treats
existing high-fidelity propagators as ``black boxes''. The two primary
estimation methods for coeffiare:

\begin{enumerate}
\def\labelenumi{\arabic{enumi}.}
\item
  \textbf{Pseudo-Spectral Projection:} Exploits orthogonality by
  projecting the model output onto each basis function using numerical
  integration (e.g., Gaussian quadrature or Smolyak sparse grids).
\item
  \textbf{Least-Squares Regression:} Minimizes the residual between the
  expansion and model evaluations at a set of sampled points (e.g.,
  Latin Hypercube Sampling).
\end{enumerate}

\begin{Shaded}
\begin{Highlighting}[]
\NormalTok{graph TD}
\NormalTok{    A[Initial PDF: p\_xi] {-}{-}\textgreater{} B[Generate Inputs: xi realization via Rosenblatt or Cholesky]}
\NormalTok{    B {-}{-}\textgreater{} C[Select Collocation Nodes via Quadrature or Sampling]}
\NormalTok{    C {-}{-}\textgreater{} D[Propagate each Node through High{-}Fidelity ODEs]}
\NormalTok{    D {-}{-}\textgreater{} E[Obtain Propagated Realizations: X\_ti]}
\NormalTok{    E {-}{-}\textgreater{} F[Solve for Spectral Coefficients: c\_j]}
\NormalTok{    F {-}{-}\textgreater{} G[Regression: Least{-}Squares Solution or Pseudo{-}Spectral Projection]}
\NormalTok{    F {-}{-}\textgreater{} H[Spectral Projection: Numerical Integration]}
\NormalTok{    G {-}{-}\textgreater{} I[Final PCE Surrogate: Analytic Response Surface]}
\NormalTok{    H {-}{-}\textgreater{} I}
\NormalTok{    I {-}{-}\textgreater{} J[Derive Moments / Sobol Indices Analytically]}
\NormalTok{    I {-}{-}\textgreater{} K[Pc Calculation: Rapid Sampling of Surrogate]}
\end{Highlighting}
\end{Shaded}

\subsection{Variations of Expansion and Surrogate
Methods}\label{variations-of-expansion-and-surrogate-methods}

To address specific limitations---such as non-standard distributions,
discontinuities, or high dimensionality---the sources identify several
key variations of the PCE method.

\begin{itemize}
\item
  \textbf{Generalized Polynomial Chaos (gPC):} Built upon the
  \textbf{Wiener-Askey scheme}, gPC matches specific polynomial families
  to standard input distributions to ensure optimal convergence. For
  example, \textbf{Hermite polynomials} are optimal for Gaussian inputs,
  while \textbf{Legendre polynomials} are used for uniform inputs.
\item
  \textbf{Arbitrary Polynomial Chaos (APC):} A \textbf{data-driven
  approach} that circumvents the need for a presumed parametric
  distribution. APC constructs an orthogonal basis using only a finite
  number of moments calculated directly from raw sampling data (e.g.,
  from an admissible region in orbit determination).
\item
  \textbf{Sparse PCE and Compressive Sensing:} To combat the curse of
  dimensionality, modern frameworks exploit the
  \textbf{sparsity-of-effects principle}---the idea that systems are
  dominated by low-order interactions. Solvers such as
  \textbf{Orthogonal Matching Pursuit (OMP)}, \textbf{Least Angle
  Regression (LARS)}, or \(l_1\)-regularized \textbf{Compressive
  Sensing} identify only the most significant coefficients, allowing for
  accurate surrogates using significantly fewer function evaluations
  than terms in the full basis.
\item
  \textbf{Multi-Element PCE (ME-PCE):} For systems with sharp gradients,
  discontinuities, or long-term integration errors, global polynomials
  often fail. ME-PCE utilizes \textbf{domain decomposition}, splitting
  the random input space into smaller sub-regions (elements) and fitting
  local PCE surrogates to each.
\item
  \textbf{Hybrid Meta-Modeling (PC-Kriging):} PCK is a hybrid
  metamodeling technique that combines the global trend-capturing
  capabilities of PCE with the local interpolation accuracy of
  \textbf{Kriging} (Gaussian Process Regression).

  \begin{itemize}
  \item
    \textbf{Formula:}
    \(M_{PCK}(\boldsymbol{X}) = \sum_{i=1}^{P} \alpha_i \Psi_i(\boldsymbol{X}) + \sigma^2 Z(\boldsymbol{X})\).
  \item
    \textbf{Methodology:} The PCE expansion serves as the trend function
    for Universal Kriging, and the Kriging part handles the local
    residuals. PCK often requires an order of magnitude fewer training
    points than standard Kriging or PCE to reach the same accuracy.
  \end{itemize}
\item
  \textbf{GMM-PCE Hybrid methods:} For cases with extreme nonlinearity
  or long propagation times, \textbf{Gaussian Mixture Models (GMMs)} are
  combined with PCE.

  \begin{itemize}
  \tightlist
  \item
    \textbf{Methodology:} The initial large Gaussian distribution is
    split into \(N\) smaller Gaussian sub-components (h-refinement). A
    lower-order PCE is then applied to each sub-component
    (p-refinement). This ``hp-refinement'' strategy reduces the domain
    of approximation for each element, allowing for rapid convergence
    where global PCE might fail.
  \end{itemize}
\item
  \textbf{Dimension Reduction (AS-NIPC \& POD-PCE):} To handle
  high-dimensional state spaces (e.g., constellation sets or
  high-fidelity fields), dimension reduction techniques are integrated:

  \begin{itemize}
  \item
    \textbf{Active Subspace (AS):} Identifies lower-dimensional
    directions in the input space that most influence the output,
    performing NIPC only in this active subspace.
  \item
    \textbf{Proper Orthogonal Decomposition (POD):} Used for ``field''
    outputs, POD identifies a latent space to reduce output
    dimensionality before mapping inputs to latent coordinates via PCE
    or PCK.
  \end{itemize}
\item
  \textbf{Multi-Fidelity (MF) Extensions:} Multi-Fidelity PCE
  incorporates information from a hierarchy of models. It uses a
  low-fidelity model (e.g., 2-body dynamics) as a foundation and a
  lower-order PCE to resolve the \textbf{model discrepancy} (additive or
  multiplicative) from a limited set of high-fidelity evaluations.
\end{itemize}

\subsection{Summary of Strengths and
Limitations}\label{summary-of-strengths-and-limitations}

{\def\LTcaptype{none} % do not increment counter
\begin{longtable}[]{@{}
  >{\raggedright\arraybackslash}p{(\linewidth - 4\tabcolsep) * \real{0.3333}}
  >{\raggedright\arraybackslash}p{(\linewidth - 4\tabcolsep) * \real{0.3333}}
  >{\raggedright\arraybackslash}p{(\linewidth - 4\tabcolsep) * \real{0.3333}}@{}}
\toprule\noalign{}
\begin{minipage}[b]{\linewidth}\raggedright
Feature
\end{minipage} & \begin{minipage}[b]{\linewidth}\raggedright
PCE Strength
\end{minipage} & \begin{minipage}[b]{\linewidth}\raggedright
Methodological Limitation
\end{minipage} \\
\midrule\noalign{}
\endhead
\bottomrule\noalign{}
\endlastfoot
\textbf{Statistical Analysis} & Moments and \textbf{Sobol Indices}
(sensitivity) are computed analytically. & Requires
square-measurable/finite variance stochastic solutions. \\
\textbf{Computational Efficiency} & Massive speed-up relative to MC; can
resolve collision risk with \textless300 evaluations vs.~\(10^5\) MC
trials. & Significantly slower than linear approximations or the
Multi-Fidelity method. \\
\textbf{Modeling Flexibility} & Non-intrusive; treats complex dynamics
as a black box. & Standard global versions struggle with discontinuities
or steep local variations. \\
\end{longtable}
}

\section{Differential Algebra (DA)}\label{differential-algebra-da}

In the larger context of \textbf{Polynomial and Expansion Methods}, the
sources define \textbf{Differential Algebra (DA)} as a computational
technique that facilitates the nonlinear propagation of uncertainties by
substituting classical real algebra with \textbf{Taylor series
expansions}. While related to other expansion techniques like Polynomial
Chaos Expansions (PCE) and State Transition Tensors (STT), DA is
distinguished by its ability to calculate higher-order derivatives
automatically.

\subsection{Mathematical Foundations of Differential Algebra
methods}\label{mathematical-foundations-of-differential-algebra-methods}

The core methodology of DA involves substituting the classical real
algebra of floating-point numbers with a truncated power series algebra,
typically denoted as \(^kD_n\), representing polynomials of order \(k\)
in \(n\) variables.

\begin{itemize}
\item
  \textbf{The Algebra \(kDn\):} This algebra consists of equivalence
  classes of functions that share the same Taylor polynomial up to order
  \(k\) in \(n\) variables centered at \(x_0\).
\item
  \textbf{State Initialization}: Instead of a single initial state
  \(\mathbf{x}_0\), the state is initialized as a DA variable
  \([\mathbf{x}] = \mathbf{x}_0 + \delta\mathbf{x}\), where
  \(\delta\mathbf{x}\) represents a deviation or uncertainty around the
  nominal reference.
\item
  \textbf{Automatic Derivation}: By overloading standard mathematical
  operators, any numerical integration scheme (e.g., Runge-Kutta) can be
  implemented to propagate the Taylor expansion of the flow
  automatically.
\item
  \textbf{Coefficient Cardinality:} The maximum number of coefficients
  \(N_{coeff}\) required to represent a polynomial of order \(k\) in
  \(n\) variables is given by: \[
  N_{coeff} = \frac{(n + k)!}{n!k!}
  \]

  For instance, a 10th-order expansion in 6 dimensions (10D6) requires
  8,008 coefficients.
\item
  \textbf{The Taylor Map}: The primary output is a high-order Taylor
  polynomial map relating initial deviations to the final state.
\end{itemize}

Mathematically, the transformation is expressed as: \[
[\mathbf{y}] = f([\mathbf{x}]) = \mathcal{T}^k_y(\delta \mathbf{x}) = \sum_{p_1 + \dots + p_n \le k} c_{p_1 \dots p_n} \cdot \delta x_1^{p_1} \dots \delta x_n^{p_n} 
\]

\textbf{Terms Definition}:

\begin{itemize}
\item
  \(\mathcal{T}^k_y\): The \(k\)-th order Taylor expansion of the
  dependent variable \(\mathbf{y}\).
\item
  \(c_{p_1 \dots p_n}\): Taylor coefficients representing high-order
  partial derivatives of the state with respect to the initial
  conditions.
\item
  \(\delta x_i\): The \(i\)-th component of the initial state deviation
  vector.
\item
  \(k\): The expansion order, which determines the accuracy of the
  nonlinear capture.
\item
  \(\mathbf{y}\): The dependent variable whose uncertainty is being
  propagated.
\end{itemize}

\subsection{Operational Methodology: High-Order Flow
Expansion}\label{operational-methodology-high-order-flow-expansion}

The DA methodology facilitates nonlinear UP by propagating a
``neighborhood'' of states rather than a single point.

\begin{enumerate}
\def\labelenumi{\arabic{enumi}.}
\item
  \textbf{DA Initialization:} The initial state \(x_0\) is initialized
  as a DA variable \([x_0] = \bar{x}_0 + \delta x_0\), where
  \(\bar{x}_0\) is the nominal mean and \(\delta x_0\) represents
  first-order variations.
\item
  \textbf{Polymorphic Integration:} The governing ordinary differential
  equations (ODEs), \(\dot{x} = f(x, t)\), are integrated using standard
  schemes (e.g., Runge-Kutta) where every real-number operation is
  replaced by its adjoint operation in the space of Taylor polynomials.
\item
  \textbf{Map Extraction:} The result is the \(k\)-th order Taylor
  expansion of the flow \(\phi(t; x_0, \alpha)\). This map is analytic
  and can be used for near-instantaneous evaluations of displaced
  initial conditions.
\item
  \textbf{Statistical Mapping:} Terminal moments are computed by
  applying the expectation operator \(\mathbb{E}[\cdot]\) to the
  resulting Taylor map.

  \begin{itemize}
  \item
    \textbf{Terminal Mean:}
    \(\mu_i(t) = \sum_{|\alpha| \le k} c_{\alpha}(t) \mathbb{E}[\delta x_0^{\alpha}]\).
  \item
    \textbf{Terminal Covariance:}
    \(P_{ij}(t) = \mathbb{E}[x_i(t)x_j(t)] - \mu_i(t)\mu_j(t)\), which
    involves algebraic operations on the coefficients and initial
    moments.
  \end{itemize}
\end{enumerate}

\subsection{Algorithmic Implementation
Flowcharts}\label{algorithmic-implementation-flowcharts}

\subsubsection{Standard DA Uncertainty
Propagation}\label{standard-da-uncertainty-propagation}

\begin{Shaded}
\begin{Highlighting}[]
\NormalTok{graph TD}
\NormalTok{    A[Initial PDF: Mean μ0, Covariance P0] {-}{-}\textgreater{} B[Initialize DA State: x\_0 = μ0 + δx0]}
\NormalTok{    B {-}{-}\textgreater{} C[Select Integration Scheme: e.g., RK45]}
\NormalTok{    C {-}{-}\textgreater{} D[Propagate via DA Arithmetic: Substitute real ops with polynomial ops]}
\NormalTok{    D {-}{-}\textgreater{} E[Obtain Terminal Taylor Map: T\_tf]}
\NormalTok{    E {-}{-}\textgreater{} F\{Analysis Type?\}}
\NormalTok{    F {-}{-} Analytic Moments {-}{-}\textgreater{} G[Evaluate E[T\_tf] and Covariance via Moment Mapping]}
\NormalTok{    F {-}{-} Monte Carlo {-}{-}\textgreater{} H[Sample δx0 and Evaluate T\_tf point{-}wise]}
\NormalTok{    G {-}{-}\textgreater{} I[Output Post{-}Propagation Statistics]}
\NormalTok{    H {-}{-}\textgreater{} I}
\end{Highlighting}
\end{Shaded}

\subsubsection{DA with Automatic Domain Splitting
(ADS)}\label{da-with-automatic-domain-splitting-ads}

\begin{Shaded}
\begin{Highlighting}[]
\NormalTok{graph TD}
\NormalTok{    A[Uncertainty Volume: Ω] {-}{-}\textgreater{} B[Propagate Taylor Map]}
\NormalTok{    B {-}{-}\textgreater{} C[Estimate Truncation Error: Order k+1 terms]}
\NormalTok{    C {-}{-}\textgreater{} D\{Error \textgreater{} Threshold?\}}
\NormalTok{    D {-}{-} Yes {-}{-}\textgreater{} E[Split Domain: Divide Ω into Ω\_left, Ω\_right]}
\NormalTok{    E {-}{-}\textgreater{} F[Propagate Patch of Local Maps]}
\NormalTok{    D {-}{-} No {-}{-}\textgreater{} G[Continue Integration]}
\NormalTok{    F {-}{-}\textgreater{} H[Reconstruct Patched Global Manifold]}
\NormalTok{    G {-}{-}\textgreater{} H}
\end{Highlighting}
\end{Shaded}

\subsection{Variations and Algorithmic
Enhancements}\label{variations-and-algorithmic-enhancements}

To overcome the limitations of local expansions and computational
growth, several advanced variants of the DA framework exist.

\begin{itemize}
\item
  \textbf{Automatic Domain Splitting (ADS/LOADS):} In highly nonlinear
  regimes, a single Taylor expansion may fail to maintain accuracy over
  the entire uncertainty volume. ADS monitors the truncation error
  (e.g., the magnitude of terms of order \(k+1\)) and adaptively
  subdivides the uncertainty domain into smaller sub-regions when error
  thresholds are exceeded.
\item
  \textbf{DA-based Monte Carlo (DAMC):} Instead of thousands of
  expensive numerical integrations, DAMC samples from the initial PDF
  and evaluates the samples using the DA-derived Taylor map. This
  typically achieves a 10-fold reduction in computation time for order
  \(k=3\).
\item
  \textbf{Hybrid GMM-DA:} The initial uncertainty is split into a
  \textbf{Gaussian Mixture Model (GMM)}, and each kernel is nonlinearly
  propagated via a high-order DA map. This approach is self-adaptive and
  avoids manual derivation of high-order derivatives.
\item
  \textbf{DA-Based Higher-Order (DAHO)}: This method analytically
  computes the moments (mean, covariance, skewness, kurtosis) of the
  transformed probability density function (PDF). This is achieved by
  computing the expectation of the monomials in the Taylor expansion.

  For a Gaussian initial distribution, the expectation of monomials is
  computed using \textbf{Isserlis' (Wick's) Formula}: \[
  E\{x_1^{s_1} \dots x_n^{s_n}\} = \begin{cases} 0 & \text{if } \sum s_i \text{ is odd} \\ Haf(\mathbf{P}) & \text{if } \sum s_i \text{ is even} \end{cases}
  \]

  \textbf{Terms Definition}:

  \begin{itemize}
  \item
    \(E\{\cdot\}\): The expectation operator.
  \item
    \(Haf(\mathbf{P})\): The Hafnian of the covariance matrix
    \(\mathbf{P}\), which is a sum over all permutations of product
    terms of the covariance components.
  \end{itemize}
\item
  \textbf{DA-LOADS Pruning:} Utilizes polynomial bounding techniques to
  discard regions of the state manifold that do not intersect with
  observations during initial orbit determination (IOD).
\item
  \textbf{Multifidelity DA (MF-DA)}:

  \begin{itemize}
  \item
    \textbf{Methodology}: This approach combines low-fidelity (LF)
    analytical models (like SGP4) propagated in DA with high-fidelity
    (HF) point-wise numerical integrations. The LF DA propagation
    provides the nonlinear ``shape'' of the uncertainty, while HF
    integrations of the polynomial centers correct the absolute position
    accuracy.
  \item
    \textbf{Efficiency}: This method can achieve speed-ups of 15--20x
    compared to full HF numerical DA integrations while maintaining high
    accuracy.
  \end{itemize}
\end{itemize}

\subsection{Differential Algebra in the Expansion
Taxonomy}\label{differential-algebra-in-the-expansion-taxonomy}

Within the broader context of Expansion and Surrogate methods, DA is
distinguished by its \textbf{intrusive nature} and \textbf{problem
independence}.

\begin{itemize}
\item
  \textbf{DA vs.~State Transition Tensors (STT):} STTs require the
  manual derivation and integration of complex high-order variational
  equations, which becomes prohibitively complex for high-fidelity
  models. DA automates this process by overloading the integrator's
  basic operations, making it essentially ODE-independent.
\item
  \textbf{DA vs.~Polynomial Chaos (PCE):} PCE is a
  \textbf{non-intrusive} ``black-box'' spectral technique that
  characterizes uncertainty using a sum of orthogonal polynomials of
  random variables. In contrast, DA focuses on the Taylor expansion of
  the state and requires the system dynamics to be continuous and
  differentiable (\(C^{k+1}\)).
\item
  \textbf{Coordinate Frame Dependency:} Unlike methods restricted to
  specific coordinates to minimize nonlinearity, DA-based methods are
  largely unaffected by coordinate selection (e.g., Cartesian
  vs.~Element space) because they can adjust expansion orders to
  maintain accuracy.
\item
  \textbf{Handling Discontinuities}: Traditional DA requires continuous
  and differentiable dynamics. However, advanced implementations (like
  Thales' PACE {[}Bignon et al. (2015){]}) approximate
  discontinuities---such as shadow entry/exit in solar radiation
  pressure---using infinitely differentiable functions like the
  arctangent to remain within the DA framework.
\item
  \textbf{The Curse of Dimensionality:} Like all expansion methods, DA
  suffers from factorial growth in terms (\(N_{coeff}\)) as order \(k\)
  increases. Modern extensions like \textbf{Distribution Transport (DT)}
  attempt to mitigate this by organizing coefficients into vectors for
  parallel computation.
\end{itemize}

\subsection{Applications and
Integration}\label{applications-and-integration}

DA is often used to enhance or complement other uncertainty
quantification frameworks:

\begin{itemize}
\item
  \textbf{Enhancing Monte Carlo (MC):} DA can significantly speed up
  Monte Carlo simulations by replacing expensive numerical integrations
  with simple \textbf{evaluations of the Taylor-expanded dynamics},
  maintaining accuracy while reducing computation time.
\item
  \textbf{Multi-Fidelity (MF) Integration:} Researchers have developed
  methods combining \textbf{MF with DA and Gaussian Mixture Models
  (GMMs)}; in these frameworks, the nonlinearity index for the
  low-fidelity step is determined using Taylor expansions computed via
  DA.
\end{itemize}

\subsection{Requirements and
Limitations}\label{requirements-and-limitations}

Despite its efficiency, DA has strict mathematical requirements that
limit its general applicability:

\begin{itemize}
\item
  \textbf{Differentiability Requirement:} The governing physical
  equations must be both \textbf{continuous and differentiable}.
\item
  \textbf{Discontinuity Barriers:} This requirement makes DA unsuitable
  for scenarios involving discontinuities, such as the step changes in
  solar radiation pressure that occur when a satellite enters or exits
  the Earth's shadow.
\end{itemize}

\section{State Transition Tensors (STT)}\label{sec-stt}

In the rigorous taxonomy of orbital uncertainty quantification (UQ),
\textbf{State Transition Tensors (STT)} represent a high-fidelity
\textbf{semi-analytic expansion method}. Situated between first-order
linear approximations like LinCov and non-intrusive surrogate models
like PCE, STTs facilitate the nonlinear propagation of uncertainties by
applying higher-order Taylor series expansions to the deviation of a
system's state from a nominal trajectory.

\subsection{Core Mechanics and Relationship to Linear
Approximation}\label{core-mechanics-and-relationship-to-linear-approximation}

STT based methods function by analytically mapping initial uncertainties
using expansions that describe nonlinear behavior within a limited
region of the state space.

\begin{itemize}
\item
  \textbf{Extension of the STM:} The sources note that a first-order
  Taylor expansion in this framework corresponds exactly to solving for
  the \textbf{State Transition Matrix (STM)} used in the
  \textbf{Linearised Covariance (LinCov)} method (See
  Section~\ref{sec-lincov}) .
\item
  \textbf{Handling Nonlinearity:} By solving for higher-order tensors,
  the method can provide significantly better results than standard
  linearisation for cases involving \textbf{strong nonlinearity}.
\item
  \textbf{Algebraic Propagation:} Once the tensors are obtained by
  numerically integrating along a nominal trajectory, they provide a
  means to compute the nonlinear propagation of the first two
  statistical moments (mean and covariance) through relatively simple
  \textbf{algebraic operations}.
\end{itemize}

\subsection{Mathematical Architecture of
STTs}\label{mathematical-architecture-of-stts}

The fundamental premise of the STT method is that the state deviation
\(\delta \mathbf{x}(t)\) at a future time can be expressed as a function
of the initial deviation \(\delta \mathbf{x}_0\) using a \(P\)-order
Taylor series expansion about a nominal trajectory \(\mathbf{x}^*(t)\).

\subsubsection{The High-Order Taylor
Map}\label{the-high-order-taylor-map}

The \(i\)-th component of the state deviation at a future time \(t\) is
approximated as: \[
\delta x_i(t) \approx \sum_{p=1}^P \frac{1}{p!} \Phi_{i,k_1 \dots k_p}(t, t_0) \delta x_{0,k_1} \dots \delta x_{0,k_p}
\]

\textbf{Term Definitions:}

\begin{itemize}
\item
  \(\delta x_i(t)\): The \(i\)-th component of the state deviation at
  time \(t\).
\item
  \(\Phi_{i,k_1 \dots k_p}\): The \(p\)-th order \textbf{State
  Transition Tensor}, defined as the \(p\)-th partial derivative of the
  solution flow with respect to initial conditions:
  \(\frac{\partial^p x_i}{\partial x_{0,k_1} \dots \partial x_{0,k_p}}\).
\item
  \(\delta x_{0,k_j}\): The \(k_j\)-th component of the initial state
  deviation.
\item
  \(P\): The truncation order of the expansion (typically \(P=2\) or
  \(3\)).
\end{itemize}

\subsubsection{Differential Equations for Tensor
Propagation}\label{differential-equations-for-tensor-propagation}

STTs are obtained by integrating a set of coupled ordinary differential
equations (ODEs) derived from the system dynamics
\(\dot{\mathbf{x}} = \mathbf{f}(\mathbf{x}, t)\). For example, the
second-order STT evolves according to: \[
\dot{\Phi}_{i,ab} = f_{i,\alpha}^\star \Phi_{\alpha,ab} + f_{i,\alpha\beta}^\star \Phi_{\alpha,a} \Phi_{\beta,b}
\]

\textbf{Term Definitions:}

\begin{itemize}
\item
  \(f_{i,\alpha}^\star, f_{i,\alpha\beta}^\star\): The first and second
  partial derivatives (Jacobian and Hessian) of the dynamics
  \(\mathbf{f}\) evaluated along the nominal trajectory.
\item
  \(\alpha, \beta\): Dummy indices using Einstein summation notation.
\item
  \(\Phi_{\alpha,a}, \Phi_{\beta,b}\): First-order STT components that
  couple into the second-order evolution.
\end{itemize}

\subsection{Methodology: Moment Propagation via Algebraic
Mapping}\label{methodology-moment-propagation-via-algebraic-mapping}

Unlike sampling methods that propagate individual points, STT-based
uncertainty propagation maps the initial statistical moments directly to
the terminal state through simple algebraic operations.

\begin{enumerate}
\def\labelenumi{\arabic{enumi}.}
\item
  \textbf{Reference Integration:} Solve the ODEs for the nominal state
  \(\mathbf{x}(t)\) and the tensors \(\phi\) up to order \(P\).
\item
  \textbf{Initial Moment Characterization:} Compute the raw moments of
  the initial PDF (e.g., using the moment-generating function for a
  Gaussian or uniform distribution).
\item
  \textbf{Statistical Mapping:} Use the Taylor expansion and the
  linearity of the expectation operator to compute the terminal mean
  (\(\hat{\mu}\)) and covariance (\(\hat{\mathbf{P}}\)): \[
  \delta \mu_i(t) \approx \sum_{p=1}^P \frac{1}{p!} \phi_{i, k_1 \dots k_p} \mathbb{E}[\delta x_{0, k_1} \dots \delta x_{0, k_p}]
  \]
\item
  \textbf{Gaussian Moment Property:} If the initial distribution is
  Gaussian, higher-order moments are analytically linked to the
  covariance via \textbf{Isserlis' Theorem} (or Wick's formula),
  allowing for precise reconstruction of non-Gaussian terminal shapes
  using only the initial mean and covariance.
\end{enumerate}

\subsection{Algorithm: The STT Uncertainty
Pipeline}\label{algorithm-the-stt-uncertainty-pipeline}

The following flowchart delineates the procedural execution of an
STT-based uncertainty propagation campaign.

\begin{Shaded}
\begin{Highlighting}[]
\NormalTok{graph TD}
\NormalTok{    A[Initial PDF: μ0, P0] {-}{-}\textgreater{} B[Generate Initial Moments: E δx0\^{}p via Isserlis]}
\NormalTok{    B {-}{-}\textgreater{} C[Set Expansion Order: P]}
\NormalTok{    C {-}{-}\textgreater{} D[Numerically Integrate: Nominal State + STM + Higher{-}Order STTs]}
\NormalTok{    D {-}{-}\textgreater{} E[Obtain Terminal Tensors at tf]}
\NormalTok{    E {-}{-}\textgreater{} F[Reconstruct Terminal Mean: μ\_tf = Σ Tensors * Initial Moments]}
\NormalTok{    F {-}{-}\textgreater{} G[Reconstruct Terminal Covariance: P\_tf = Σ f Tensors, P0]}
\NormalTok{    G {-}{-}\textgreater{} H\{Capture Non{-}Gaussianity?\}}
\NormalTok{    H {-}{-} Yes {-}{-}\textgreater{} I[Reconstruct Higher Moments: Skewness \& Kurtosis]}
\NormalTok{    H {-}{-} No {-}{-}\textgreater{} J[Output Gaussian Approximation]}
\NormalTok{    I {-}{-}\textgreater{} K[Pc Calculation for Conjunction Assessment]}
\NormalTok{    J {-}{-}\textgreater{} K}
\end{Highlighting}
\end{Shaded}

\subsection{Variations and Algorithmic
Enhancements}\label{variations-and-algorithmic-enhancements-1}

To overcome the ``curse of dimensionality'' and the high computational
cost of calculating factorial-order partial derivatives, several
advanced STT variants have been developed.

\begin{itemize}
\item
  \textbf{Directional State Transition Tensors (DSTT):} Proposed to
  reduce complexity by only capturing nonlinear effects along
  ``sensitive'' directions (e.g., the maximum stretch direction of the
  Cauchy-Green Tensor). This reduces the number of integrated variables
  from \(O(n^{P+1})\) to \(O(n^2 + nm^{P-1})\), where \(m\) is the
  number of sensitive directions.
\item
  \textbf{Time-varying Directional STT (TDSTT):} Improves upon the DSTT
  by integrating the eigenvalue-eigenvector pairs over time, allowing
  uncertainty analysis at any historical point rather than just a
  predefined final epoch.
\item
  \textbf{Distribution Transport (DT):} Organizes polynomial
  coefficients into a vector form to facilitate parallel calculation
  using linear algebra libraries and integral correction methods,
  achieving speed-ups of up to 60 times over conventional STT
  implementations.
\item
  \textbf{Reduced STT (RSTT):} Assumes two-body dynamics dominate and
  retains only the secular terms to decrease the integration burden for
  high-order approximations.
\item
  \textbf{GMM-STT Hybrid Method:} Combines Gaussian Mixture Models with
  STTs. The initial uncertainty is split into small Gaussian components,
  each of which is propagated linearly or with low-order STTs to
  maintain validity over long durations.
\item
  \textbf{Non-Gaussian STT Propagation:} Extends the method to handle
  arbitrary initial PDFs (e.g., uniform distributions for asteroid
  gravity fields) by leveraging Moment Generating Functions (MGFs) to
  calculate initial expectations.
\end{itemize}

\subsection{STT in the Spectrum of Surrogate
Methods}\label{stt-in-the-spectrum-of-surrogate-methods}

Within the broader category of surrogate methods, STTs are defined by
distinct trade-offs regarding differentiability and dimensionality.

{\def\LTcaptype{none} % do not increment counter
\begin{longtable}[]{@{}
  >{\raggedright\arraybackslash}p{(\linewidth - 6\tabcolsep) * \real{0.2500}}
  >{\raggedright\arraybackslash}p{(\linewidth - 6\tabcolsep) * \real{0.2500}}
  >{\raggedright\arraybackslash}p{(\linewidth - 6\tabcolsep) * \real{0.2500}}
  >{\raggedright\arraybackslash}p{(\linewidth - 6\tabcolsep) * \real{0.2500}}@{}}
\toprule\noalign{}
\begin{minipage}[b]{\linewidth}\raggedright
Feature
\end{minipage} & \begin{minipage}[b]{\linewidth}\raggedright
State Transition Tensors (STT)
\end{minipage} & \begin{minipage}[b]{\linewidth}\raggedright
Differential Algebra (DA)
\end{minipage} & \begin{minipage}[b]{\linewidth}\raggedright
Polynomial Chaos (PCE)
\end{minipage} \\
\midrule\noalign{}
\endhead
\bottomrule\noalign{}
\endlastfoot
\textbf{Expansion Type} & Semi-Analytic Taylor Series & Automated Taylor
Polynomials & Non-Intrusive Spectral Basis \\
\textbf{Dynamics Requirement} & Continuous \& Differentiable &
Continuous \& Differentiable & Black-Box (Non-Differentiable) \\
\textbf{Input Dependence} & Distribution Agnostic & Distribution
Agnostic & Distribution Matched (Wiener-Askey) \\
\textbf{Complexity Growth} & Exponential/Factorial with Order &
Controlled via Automated Eval & Factorial (Curse of Dimensionality) \\
\textbf{Main Advantage} & Analytic moment propagation & Automated
derivative extraction & Exponential convergence for smooth systems \\
\end{longtable}
}

The primary takeaway is that STTs provide a \textbf{local higher-order
approximation} of the flow manifold. While they accurately capture the
``banana-shaped'' warping caused by \(1/r^2\) gravitational
nonlinearities, they are strictly limited to the \textbf{radius of
convergence} of the Taylor expansion. If the initial uncertainty is too
large, the expansion fails, necessitating hybrid methods such as
\textbf{GMM-STT}, which partition the state space into smaller, locally
linear Gaussian kernels.

\subsection{Performance and
Requirements}\label{performance-and-requirements}

While STT is characterized as highly efficient for evaluating the final
uncertainty, it faces strict mathematical and implementation barriers:

\begin{itemize}
\item
  \textbf{Differentiability:} Similar to DA and LinCov, STT requires the
  underlying physical dynamics of the system to be \textbf{continuous
  and differentiable}.
\item
  \textbf{Complexity Trade-off:} Although the algebraic evaluation of
  the final state is fast, the method is described as
  \textbf{computationally intensive} during the initial phase because of
  the difficulty involved in deriving and integrating the high-order
  tensors.
\end{itemize}

\subsection{Role in Conjunction Assessment
(CA)}\label{role-in-conjunction-assessment-ca-1}

In the field of orbital mechanics, the STT method is frequently combined
with the other uncertainty propagation architectures. The sources
specifically mention the development of hybrid frameworks that pair
\textbf{Gaussian Mixture Models (GMMs) with STT} to efficiently
propagate non-Gaussian uncertainties for conjunction analysis. This
allows the GMM to break a large non-Gaussian problem into smaller
Gaussian ``kernels,'' each of which is then non-linearly propagated
using the semi-analytic STT method.

\section{Multi-Fidelity (MF) Method of Uncertainty
Propagation}\label{multi-fidelity-mf-method-of-uncertainty-propagation}

In the hierarchy of uncertainty propagation (UP) methodologies, the
\textbf{Multi-Fidelity (MF) Method} represents a data-adaptive,
\textbf{efficiency-optimized} paradigm designed to reconcile the
computational burden of high-fidelity (HF) orbital dynamics (meaning
force-models) with the large-sample requirements of robust statistical
characterization. MF methods bridge the gap between ``weak''
approximation sampling (Monte Carlo) and ``strong'' spectral expansions
(Polynomial Chaos), utilizing a hierarchy of dynamical models to
construct a surrogate of the propagated probability density function
(PDF) at a fraction of the cost.

\subsection{Theoretical Framework and
Assumptions}\label{sec-mf-assumptions}

The fundamental premise of MF methods is that a lower-fidelity model can
provide a sufficient basis for the space of propagated samples, which is
then corrected by a small subset of higher-fidelity realizations. This
approach is particularly effective in astrodynamics because the ``main
problem'' (two-body dynamics plus \(J_2\) perturbations) captures the
bulk of orbital motion, while HF perturbations (higher-order gravity,
atmospheric drag, solar radiation pressure) act as refinements.

\subsubsection{Assumptions and
Definitions}\label{assumptions-and-definitions}

\begin{itemize}
\item
  \textbf{Quantity of Interest (QoI):} These are the variables of
  interest whose uncertainties are needed to be propagated through the
  dynamics, such as the orbital state vector \(x \in \mathbb{R}^6\).
\item
  \textbf{Model Hierarchy:} In astrodynamics, the LF model typically
  employs two-body Keplerian dynamics or \(J_2\) perturbations, while
  the HF model incorporates high-degree gravity fields, atmospheric
  drag, and solar radiation pressure.
\item
  \textbf{Surrogate Basis:} A large ensemble of \(m\) LF samples (say a
  sample \(\mathbf{x}\) is a vector \(\in \mathbb{R}^n\)) defines a
  basis for the space of propagated quantities of interest (QoI). From
  this, a small subset \(r\) of ``important samples'' (where
  \(r < n \ll m\)) is identified for HF evaluation. `
\item
  \textbf{Correction Mapping:} The information gained from the \(r\) HF
  evaluations is used to ``lift'' the remaining \(m-r\) LF samples
  toward the HF truth using a stochastic collocation surrogate.
\end{itemize}

\subsection{Mathematical Architecture}\label{mathematical-architecture}

\subsubsection{Method 1: Stochastic Collocation (SC)
Multi-Fidelity}\label{sec-mf-sc}

Developed by Jones and Weisman (2019), this non-intrusive method uses a
large ensemble of LF samples to identify a subset of ``important''
samples for HF propagation.

\begin{itemize}
\item
  \textbf{Initial Ensemble:} A set of \(m\) random inputs
  \(\Xi = \{\xi_i\}_{i=1}^m\) is generated from the initial probability
  density function (PDF).
\item
  \textbf{LF Basis Construction:} The entire ensemble is propagated
  using an LF model (e.g., SGP4 or two-body dynamics) to define a
  subspace \(X_L(\Xi) \subseteq \mathcal{X}\).
\item
  \textbf{Important Sample Selection:} Identifying the \(r\) nodes is an
  optimization problem solved via a \textbf{greedy algorithm} that
  maximizes the distance between the next selected node and the existing
  basis. This is implemented by solving the pivoted Cholesky
  decomposition of the Gramian matrix:
  \begin{equation}\protect\phantomsection\label{eq-mf-pivoted-cholesky-decomp}{
  [X_L]^T G_L [X_L] = A^T C_h C_h^T A
  }\end{equation}

  \textbf{Terms defined:}

  \begin{itemize}
  \item
    \(X_L\): Matrix of LF propagated samples.
  \item
    \(G_L\): Gramian matrix, where
    \([G_L]_{i,j} = x_L(\xi_i) \cdot x_L(\xi_j)\) (\(x_L(\xi)\) is the
    LF propagated state vector for sample \(\xi\)).
  \item
    \(A\): Permutation/pivot matrix that identifies the indices of the
    important samples.
  \item
    \(C_h\): Lower triangular Cholesky factor.
  \end{itemize}

  See Section~\ref{sec-appendix-pivoted-cholesky} for a detailed
  derivation of the pivoted Cholesky algorithm and its application to MF
  uncertainty propagation.
\item
  \textbf{\(c_\ell(\xi)\) Coefficient Computation:} The expansion
  coefficients \(c_\ell(\xi)\) in Equation~\ref{eq-mf-surrogate} are
  computed by solving the linear system: \[
   \mathbf{C_h C_h}^T \mathbf{c}(\mathbf{\xi}) = \mathbf{g}
   \]

  where \(g_i = x_L(\xi) \cdot x_L(\xi_i)\) for \(i = 1, \ldots, r\).

  These coefficients effectively map the LF realization of a general
  sample \(\xi\) to the HF space defined by the \(r\) important samples
  via Equation~\ref{eq-mf-surrogate}.
\item
  \textbf{HF Correction:} The \(r\) samples are propagated through the
  HF model, and a surrogate is built using the LF-derived coefficients:
  \begin{equation}\protect\phantomsection\label{eq-mf-surrogate}{
  \hat{x}_H(\xi) = \sum_{l=1}^r c_\ell(\xi)x_H(\xi_\ell^*)
  }\end{equation}

  \textbf{Terms defined:}

  \begin{itemize}
  \item
    \(\hat{x}_H(\xi)\): The approximated HF state vector for a given
    input \(\xi\).
  \item
    \(x_H(\xi_\ell)\): The exact HF state of the \(\ell\)-th ``important
    sample''.
  \item
    \(c_\ell(\xi)\): Expansion coefficients derived from the LF model to
    relate the general sample \(\xi\) to the important nodes.
  \item
    \(r\): The rank of the surrogate (number of HF evaluations).
  \end{itemize}
\end{itemize}

Equation~\ref{eq-mf-surrogate} is what is referred to as the
``low-rank'' \textbf{stochastic collocation surrogate}. The coefficients
\(c_\ell(\xi)\) are computed by solving a linear system derived from the
LF model, ensuring that the surrogate accurately captures the HF
dynamics while minimizing the number of expensive HF evaluations.

\paragraph{\texorpdfstring{\textbf{Algorithmic
Flowchart}}{Algorithmic Flowchart}}\label{algorithmic-flowchart}

\begin{Shaded}
\begin{Highlighting}[]
\NormalTok{graph TD}
\NormalTok{    A[Start: Initial PDF p(x0)] {-}{-}\textgreater{} B[Generate random inputs ensemble Xi with size m]}
\NormalTok{    B {-}{-}\textgreater{} C[Propagate ensemble to tf using LF model]}
\NormalTok{    C {-}{-}\textgreater{} D[Construct matrix of LF samples XL]}
\NormalTok{    D {-}{-}\textgreater{} E[Compute Pivoted Cholesky Decomposition of XL Gramian]}
\NormalTok{    E {-}{-}\textgreater{} F[Select r important samples Xi* from pivot matrix A]}
\NormalTok{    F {-}{-}\textgreater{} G[Propagate only Xi* samples to tf using HF model]}
\NormalTok{    G {-}{-}\textgreater{} H[Compute surrogate coefficients c\_l using LF basis L]}
\NormalTok{    H {-}{-}\textgreater{} I[Reconstruct/Correct entire HF ensemble using c\_l and HF realizations]}
\NormalTok{    I {-}{-}\textgreater{} J[Result: Refined Propagated PDF p(xf)]}
\NormalTok{    J {-}{-}\textgreater{} K[Stop]}
\end{Highlighting}
\end{Shaded}

\subsubsection{Method 2: Differential Algebra (DA) and Adaptive GMM
Multi-Fidelity}\label{method-2-differential-algebra-da-and-adaptive-gmm-multi-fidelity}

Developed by Fossà et al., this method combines the local accuracy of
Taylor expansions with the global flexibility of Gaussian Mixture Models
(GMMs).

\begin{itemize}
\item
  \textbf{Initialization:} The initial PDF is represented as a GMM where
  each kernel is initialized in the DA framework as a Taylor polynomial.
\item
  \textbf{Adaptive Splitting (LOADS):} The Low-Order Automatic Domain
  Splitting (LOADS) algorithm monitors a DA-based nonlinearity index
  (\(NLI\)) during LF propagation. If the \(NLI\) exceeds a threshold
  \(\epsilon_\nu\), the polynomial is split into child kernels to
  maintain quasi-linearity.
\item
  \textbf{LF Propagation:} Statistics are mapped through the LF dynamics
  (e.g., SGP4) via polynomial evaluation.
\item
  \textbf{HF Rectification:} The constant part (center) of each Taylor
  expansion is propagated point-wise in HF dynamics. The accuracy is
  restored by re-centering the LF nilpotent part around the HF
  trajectory: \[
  [x_{MF}^{(l)}(t_f)] = \mu_{HF}^{(l)}(t_f) + \{[x_{LF}^{(l)}(t_f)] - \bar{x}_{LF}^{(l)}(t_f)\}
  \]

  where \([x_{MF}^{(l)}]\) is the MF polynomial, \(\mu_{HF}^{(l)}\) is
  the HF-propagated mean, and the term in braces is the LF nilpotent
  (varying) part.
\end{itemize}

\paragraph{\texorpdfstring{\textbf{Algorithmic
Flowchart}}{Algorithmic Flowchart}}\label{algorithmic-flowchart-1}

The following flowchart illustrates the bi-fidelity uncertianty
propagation using Differential Algebra and Adaptive GMM as described in
the sources:

\begin{Shaded}
\begin{Highlighting}[]
\NormalTok{graph TD}
\NormalTok{    A[Start: Initial PDF p(x0)] {-}{-}\textgreater{} B[Initialize DA variables or Particles]}
\NormalTok{    B {-}{-}\textgreater{} C[LF Propagation: Map uncertainty through cheap model]}
\NormalTok{    C {-}{-}\textgreater{} D\{Nonlinearity Check?\}}
\NormalTok{    D {-}{-} Yes {-}{-}\textgreater{} E[LOADS/AEGIS: Split domains/kernels]}
\NormalTok{    E {-}{-}\textgreater{} C}
\NormalTok{    D {-}{-} No {-}{-}\textgreater{} F[Identify Important Samples/Kernel Means]}
\NormalTok{    F {-}{-}\textgreater{} G[HF Propagation: Point{-}wise propagation of select points]}
\NormalTok{    G {-}{-}\textgreater{} H[Polynomial/Surrogate Correction: Re{-}center or Re{-}weight]}
\NormalTok{    H {-}{-}\textgreater{} I[Result: Refined Propagated PDF p(xf)]}
\NormalTok{    I {-}{-}\textgreater{} J[Stop]}
\end{Highlighting}
\end{Shaded}

\subsection{Hybrid and Efficiency-Optimized
Frameworks}\label{hybrid-and-efficiency-optimized-frameworks}

The sources highlight several advanced hybridizations and methodological
variations that leverage MF's scalability:

\begin{itemize}
\item
  \textbf{GMM-MF (Gaussian Mixture Multi-Fidelity):} (from Jones and
  Weisman (2019)) The GMM-MF framework combines the \textbf{Gaussian
  Mixture Model (GMM)} with \textbf{stochastic collocation MF method} to
  handle non-Gaussian distributions. The GMM decomposes the initial
  uncertainty into multiple Gaussian kernels. Then sigma points
  (\(\zeta\)) are deterministically generated for each kernel, which
  correspond to random inputs \(\xi\) that are then propagated through
  multi-fidelity stochastic collocation. This allows for accurate
  propagation of non-Gaussian distributions while maintaining
  computational efficiency.
\item
  \textbf{Tri-Fidelity Models:} (from Jones and Weisman (2019)) This
  extension introduces a medium-fidelity (MF) model. LF is used to
  identify important samples, MF is used to calculate the expansion
  coefficients \(c\), and HF provides the final collocation points,
  further optimizing the accuracy-to-cost ratio.
\item
  \textbf{Stochastic Acceleration Propagation (PLASMA):} (from Fossà
  (2024)) To account for poorly modeled forces (process noise), the
  PLASMA algorithm propagates polynomial approximations of the moments
  of the PDF through stochastic differential equations (SDEs). In an MF
  framework, the initial uncertainty is handled by GMM-DA, while the
  process noise moments are propagated independently and merged \emph{a
  posteriori}.
\item
  \textbf{Analytical vs.~Semianalytical Coupling:}

  \begin{itemize}
  \item
    \emph{Analytical LF:} Uses closed-form solutions like SGP4 or
    \(J_2\)-perturbed Lagrangian coefficients for near-instantaneous
    mapping.
  \item
    \emph{Semianalytical LF:} Uses mean element theories (e.g., DSST,
    STELA) where short-period variations are removed. This allows larger
    integration steps while maintaining better fidelity than purely
    analytical models.
  \end{itemize}
\item
  \textbf{Credibilistic Filtering:} (from Jones et al. (2019)) Because
  the surrogate introduces systematic (epistemic) error rather than
  random error, recent frameworks use \textbf{Outer Probability Measures
  (OPMs)} and possibility functions \(f(x)\) to bound these errors. This
  provides a ``less over-descriptive'' representation of uncertainty for
  establishing custody of newly detected objects.
\end{itemize}

\subsection{Multi-Fidelity in Orbit
Propagation}\label{multi-fidelity-in-orbit-propagation}

A critical nuance in multi-fidelity orbit propagation is the
\textbf{state trajectory expansion}.

\begin{itemize}
\item
  \textbf{The n=6 Limitation:} Standard MF surrogates are bounded by the
  state dimension (\(n=6\)). To use more nodes (\(r > 6\)) for higher
  accuracy, Jones and Weisman (2019) expand the Quantity of Interest
  (QoI) vector to include the \textbf{state trajectory}.
\item
  \textbf{Trajectory Augmentation:} By defining the QoI as the state at
  \(T\) different time instances, the effective dimensionality becomes
  \(6T\), allowing for a significantly higher surrogate rank and better
  representation of the nonlinear manifold.
\end{itemize}

\subsection{Performance and Conjunction Assessment
Significance}\label{performance-and-conjunction-assessment-significance}

\begin{itemize}
\item
  \textbf{Computational Speedup:} MF achieves a speedup of up to
  \textbf{four orders of magnitude} (\(10^4\)) compared to brute-force
  Monte Carlo for high-fidelity LEO propagation.
\item
  \textbf{Handling Non-Gaussianity:} Unlike linear methods, MF captures
  ``banana-shaped'' warpings and asymmetric tails. In conjunction
  assessment, as few as \textbf{8 important samples} allow MF to produce
  collision probabilities (\(P_C\)) that fall within the 95\% confidence
  interval of a full MC simulation.
\item
  \textbf{Dimensionality custom:} To avoid limits on \(r\) based on
  state size (\(n=6\)), researchers use the \textbf{state trajectory}
  (state at multiple time points) to increase the effective dimension of
  the Quantity of Interest (QoI) vector, enabling a more accurate basis.
\end{itemize}

\section{Taylor Map Diffusion Method for Closed-Form solution to
Fokker-Planck
equation}\label{taylor-map-diffusion-method-for-closed-form-solution-to-fokker-planck-equation}

\subsection{Fokker-Planck Equation in Orbital
Mechanics}\label{fokker-planck-equation-in-orbital-mechanics}

The \textbf{Fokker-Planck Equation (FPE)} is a partial differential
equation that governs the time evolution of the probability density
function (PDF) of a dynamical system subject to both deterministic
forces and stochastic process noise. In the context of orbital mechanics
and conjunction assessment, the FPE is the fundamental tool for
transitioning from a point-estimate of a spacecraft's position to a
complete statistical characterization of its possible states.

The FPE describes how a probability cloud ``flows'' through state space
while simultaneously ``spreading'' due to unmodeled perturbations.

\begin{itemize}
\item
  \textbf{Advection-Diffusion Duality:} The equation is composed of two
  primary operators: the \textbf{drift (advection)} term, which moves
  the PDF along deterministic trajectories (like Keplerian motion), and
  the \textbf{diffusion} term, which accounts for stochastic forcing
  such as atmospheric drag fluctuations or thrust errors.
\item
  \textbf{The Governing Equation:} For a system
  \(dX = f(X)dt + \sigma dW_t\), the PDF \(p(x, t)\) evolves as: \[
  \frac{\partial p}{\partial t} = -\nabla \cdot (f p) + \frac{1}{2} \text{Tr}(\Sigma_0 \nabla^2 p)
  \]
\end{itemize}

Where \(f\) is the deterministic drift and \(\Sigma_0\) is the diffusion
tensor.

\begin{itemize}
\tightlist
\item
  \textbf{The Curse of Dimensionality:} Traditional numerical solvers
  for the FPE rely on spatial grids. In a 6D orbital state, even a
  coarse grid of 50 points per axis requires approximately \(10^{10}\)
  evaluations per time step, making grid-based FPE solutions
  computationally intractable for real-time space operations.
\end{itemize}

By providing a closed-form solution to the FPE, TMD addresses the
long-standing ``process noise gap'' in non-intrusive orbital uncertainty
propagation.

\subsection{Theoretical Mechanics of the TMD
Framework}\label{theoretical-mechanics-of-the-tmd-framework}

The TMD method is built upon the \textbf{Quadratic Form Conservation
Theorem}, which establishes that if a distribution starts with a
specific structural form, it will maintain that form throughout its
evolution.

\begin{itemize}
\item
  \textbf{The Structural Ansatz:} The PDF is represented as
  \(p(x, t) = N(t) \mu(x, t) \exp(-F(x, t)^T Q(t) F(x, t))\). This
  formulation separates the distribution into three primary components:

  \begin{itemize}
  \item
    \textbf{\(F(x, t)\):} A nonlinear Taylor map that captures the
    \textbf{shape} of the distribution (curvatures, ``banana-shapes,''
    and asymmetries).
  \item
    \textbf{\(Q(t)\):} A precision matrix (inverse of covariance) that
    captures the \textbf{concentration} and stochastic broadening.
  \item
    \textbf{\(\mu(x, t)\):} A scalar field accounting for the
    \textbf{divergence} of non-conservative forces like drag.
  \end{itemize}
\item
  \textbf{Augmented ODE System:} Instead of solving a high-dimensional
  PDE, TMD reduces the problem to a compact system of coupled Ordinary
  Differential Equations (ODEs). For a 6D state, this involves only
  \textbf{279 scalar ODEs}, which can be solved in less than a second on
  standard hardware.
\item
  \textbf{Automatic Differentiation (AD):} The implementation leverages
  AD (e.g., via JAX) to extract the necessary Jacobian and Hessian
  information from the orbital dynamics, eliminating the need for manual
  derivation of complex 6D force models.
\end{itemize}

\subsection{Algorithm: The Taylor Diffusion
Propagator}\label{algorithm-the-taylor-diffusion-propagator}

\begin{Shaded}
\begin{Highlighting}[]
\NormalTok{graph TD}
\NormalTok{    A[Initial State X0, Covariance P0, Process Noise Σ0] {-}{-}\textgreater{} B[Initialize: Trajectory x0, Map J, Hessian H, Precision Q]}
\NormalTok{    B {-}{-}\textgreater{} C[Compute Precision Matrix Derivative ¤Q]}
\NormalTok{    C {-}{-}\textgreater{} D[Extract Map Evolution Derivatives via Automatic Differentiation]}
\NormalTok{    D {-}{-}\textgreater{} E[Identify Taylor Coefficients: V0, ¤J, ¤H]}
\NormalTok{    E {-}{-}\textgreater{} F[Apply Zero{-}Tracking Condition: ¤x0 = {-}J⁻¹V0]}
\NormalTok{    F {-}{-}\textgreater{} G[Integrate Augmented ODE System via RK4]}
\NormalTok{    G {-}{-}\textgreater{} H\{Final Time Reached?\}}
\NormalTok{    H {-}{-} No {-}{-}\textgreater{} C}
\NormalTok{    H {-}{-} Yes {-}{-}\textgreater{} I[Output Analytic PDF: \{x0, J, H, Q\}]}
\NormalTok{    I {-}{-}\textgreater{} J[Instantaneous Evaluation of Tail Probabilities]}
\end{Highlighting}
\end{Shaded}

\section{Appendix}\label{appendix}

\section{Appendix A: Pivoted Cholesky Decomposition for Multi-Fidelity
Node Selection}\label{sec-appendix-pivoted-cholesky}

In the architecture of \textbf{Multi-Fidelity (MF) Uncertainty
Propagation}, the \textbf{pivoted Cholesky decomposition} serves as the
mathematical engine for the ``greedy'' selection of \textbf{important
samples} (nodes). This approach identifies a small subset of points from
a large low-fidelity (LF) ensemble that are most representative of the
state space, ensuring that high-fidelity (HF) evaluations are spent
where they provide the most information.

\subsection{The Objective: Information
Maximization}\label{the-objective-information-maximization}

The goal of node selection is to find a set of random inputs
\(\underline{\Xi} = \{\underline{\xi}_1, \dots, \underline{\xi}_r\}\)
such that the span of their low-fidelity realizations,
\(\mathcal{X}_L(\underline{\Xi})\), provides the best possible
approximation of the full ensemble \(X_L(\Xi)\).

\begin{itemize}
\item
  \textbf{Greedy Selection Logic:} The algorithm iteratively picks the
  next node \(\xi_k\) by maximizing the distance between that sample and
  the space spanned by the previously selected nodes.
\item
  \textbf{Intractability:} Solving this optimization problem directly
  via
  \(\text{arg max dist}(\mathbf{x}(\xi), \mathcal{X}(\underline{\Xi}_{k-1}))\)
  is computationally expensive (intractable) for high-dimensional state
  spaces.
\item
  \textbf{The Cholesky Proxy:} The pivoted Cholesky decomposition
  provides an efficient linear algebra substitute to this geometric
  search by operating on the \textbf{Gramian matrix} of the low-fidelity
  samples.
\end{itemize}

\subsection{Mathematical
Architecture}\label{mathematical-architecture-1}

The method leverages the relationship between the inner products of the
propagated samples to identify the ``rank'' or ``important'' directions
of the uncertainty volume.

\subsubsection{\texorpdfstring{\textbf{The Decomposition
Formula}}{The Decomposition Formula}}\label{the-decomposition-formula}

The decomposition is defined as: \[
[X_L]^T G_L [X_L] = A^T C_h C_h^T A
\]

\textbf{Terms defined:}

\begin{itemize}
\item
  \(X_L\): The matrix of low-fidelity propagated samples
  (\(n \times m\)).
\item
  \(G_L\): The Gramian matrix where
  \([G_L]_{i,j} = x_L(\xi_i) \cdot x_L(\xi_j)\). It captures the
  similarity/redundancy between all sample pairs.
\item
  \(A\): A \textbf{pivot matrix} (permutation matrix) that reorders the
  samples based on their importance.
\item
  \(C_h\): A lower triangular Cholesky factor (\(m \times r\)).
\end{itemize}

\subsubsection{\texorpdfstring{\textbf{Weight Update and
Orthogonalization}}{Weight Update and Orthogonalization}}\label{weight-update-and-orthogonalization}

During the decomposition, the algorithm maintains a vector of
``weights'' \(w_\ell\) representing the norm of each sample. At each
step:

\begin{enumerate}
\def\labelenumi{\arabic{enumi}.}
\item
  The sample with the largest remaining weight (the most unique point)
  is selected and ``pivoted'' to the front.
\item
  The weights of all other samples are updated by subtracting their
  projection onto the newly selected node: \[
  w_t = w_t - C_{h(t,n)}^2
  \]

  This step effectively ``deflates'' the importance of samples that are
  redundant with the current basis.
\end{enumerate}

\subsection{Algorithm: Node Selection via Pivoted
Cholesky}\label{algorithm-node-selection-via-pivoted-cholesky}

The following flowchart details the procedural logic for identifying
\(r\) important samples from \(N\) low-fidelity evaluations:

\begin{Shaded}
\begin{Highlighting}[]
\NormalTok{graph TD}
\NormalTok{    A[Matrix of LF samples: X\_L] {-}{-}\textgreater{} B[Initialize Weights w: Square of norms of columns in X\_L]}
\NormalTok{    B {-}{-}\textgreater{} C[Set current rank n = 1]}
\NormalTok{    C {-}{-}\textgreater{} D[Find index i with max weight w\_i]}
\NormalTok{    D {-}{-}\textgreater{} E[Is w\_i \textless{} Tolerance?]}
\NormalTok{    E {-}{-} Yes {-}{-}\textgreater{} F[Stop: System rank reached]}
\NormalTok{    E {-}{-} No {-}{-}\textgreater{} G[Swap column n with column i in X\_L and weight vector w]}
\NormalTok{    G {-}{-}\textgreater{} H[Compute diagonal element: Ch\_n,n = sqrt w\_n]}
\NormalTok{    H {-}{-}\textgreater{} I[Compute Ch\_t,n for remaining samples t \textgreater{} n]}
\NormalTok{    I {-}{-}\textgreater{} J[Update remaining weights: w\_t = w\_t {-} Ch\_t,n\^{}2]}
\NormalTok{    J {-}{-}\textgreater{} K[n = n + 1]}
\NormalTok{    K {-}{-}\textgreater{} L\{n \textless{}= r?\}}
\NormalTok{    L {-}{-} Yes {-}{-}\textgreater{} D}
\NormalTok{    L {-}{-} No {-}{-}\textgreater{} M[Output indices of r Important Samples]}
\end{Highlighting}
\end{Shaded}

\protect\phantomsection\label{refs}
\begin{CSLReferences}{1}{1}
\bibitem[\citeproctext]{ref-acciarini:10.2514ux2f1.G008717}
Acciarini, Giacomo, Nicola Baresi, David J. B. Lloyd, and Dario Izzo.
2025. {``Nonlinear Propagation of Non-Gaussian Uncertainties.''}
\emph{Journal of Guidance, Control, and Dynamics} 48 (4): 903--13.
\url{https://doi.org/10.2514/1.G008717}.

\bibitem[\citeproctext]{ref-Adurthi2012TheCUT}
Adurthi, Nagavenkat, Puneet Singla, and Tarunraj Singh. 2012. {``The
Conjugate Unscented Transform --- an Approach to Evaluate
Multi-Dimensional Expectation Integrals.''} \emph{2012 American Control
Conference (ACC)}, 5556--61.
\url{https://doi.org/10.1109/ACC.2012.6314970}.

\bibitem[\citeproctext]{ref-bignon2015accurate}
Bignon, Emmanuel, P Mercier, Vincent Azzopardi, and R Pinéde. 2015.
{``Accurate Numerical Orbit Propagation Using Polynomial Algebra
Computational Engine Pace.''} \emph{ISSFD 2015 Congress}.

\bibitem[\citeproctext]{ref-bosorbit}
Bos, Justin. 2025. \emph{Orbit Uncertainty Propagation in Space
Situational Awareness Applied to Conjunction Assessment}.

\bibitem[\citeproctext]{ref-DETOMASO20257226}
Detomaso, Roberto, Andrea Muciaccia, and Camilla Colombo. 2025.
{``Hybrid Gaussian Mixture Model and Unscented Transformation Algorithm
for Uncertainty Propagation Within the PUZZLE Software.''}
\emph{Advances in Space Research} 75 (10): 7226--41.
https://doi.org/\url{https://doi.org/10.1016/j.asr.2025.03.026}.

\bibitem[\citeproctext]{ref-fossa2024propagation}
Fossà, Alberto. 2024. {``Propagation Multi-Fid{é}lit{é} d'incertitude
Orbitale En Pr{é}sence d'acc{é}l{é}rations Stochastiques.''} PhD thesis,
The University of Texas at Austin.

\bibitem[\citeproctext]{ref-JONES2019406}
Jones, Brandon A., and Ryan Weisman. 2019. {``Multi-Fidelity Orbit
Uncertainty Propagation.''} \emph{Acta Astronautica} 155: 406--17.
https://doi.org/\url{https://doi.org/10.1016/j.actaastro.2018.10.023}.

\bibitem[\citeproctext]{ref-2019amosconfE14J}
Jones, Brandon, Emmanuel Delande, Enrico Zucchelli, and Moriba Jah.
2019. {``{Multi-Fidelity Orbit Uncertainty Propagation with Systematic
Errors}.''} In \emph{Advanced Maui Optical and Space Surveillance
Technologies Conference}, edited by S. Ryan.

\bibitem[\citeproctext]{ref-Julier1997NewEO}
Julier, Simon J., and Jeffrey K. Uhlmann. 1997a. {``New Extension of the
Kalman Filter to Nonlinear Systems.''} \emph{Defense, Security, and
Sensing}. \url{https://api.semanticscholar.org/CorpusID:7937456}.

\bibitem[\citeproctext]{ref-Julier1997NewUKF}
Julier, Simon J., and Jeffrey K. Uhlmann. 1997b. {``{New extension of
the Kalman filter to nonlinear systems}.''} In \emph{Signal Processing,
Sensor Fusion, and Target Recognition VI}, edited by Ivan Kadar, vol.
3068. International Society for Optics; Photonics; SPIE.
\url{https://doi.org/10.1117/12.280797}.

\bibitem[\citeproctext]{ref-park-scheeres:10.2514ux2f1.20177}
Park, Ryan S., and Daniel J. Scheeres. 2006. {``Nonlinear Mapping of
Gaussian Statistics: Theory and Applications to Spacecraft Trajectory
Design.''} \emph{Journal of Guidance, Control, and Dynamics} 29 (6):
1367--75. \url{https://doi.org/10.2514/1.20177}.

\bibitem[\citeproctext]{ref-thesis_suyog}
Soman, Suyog. 2025. \emph{Orbital Uncertainty Propagation for Collision
Probability Computation}.

\bibitem[\citeproctext]{ref-sun2026recent}
Sun, Pan, Kai Cao, and Shuang Li. 2026. {``Recent Advances of Practical
Techniques for Nonlinear Uncertainty Propagation in Orbital Dynamics.''}
\emph{Science China Physics, Mechanics \& Astronomy} 69 (1): 214501.

\bibitem[\citeproctext]{ref-Tenne2003TheHOUF}
Tenne, D., and T. Singh. 2003. {``The Higher Order Unscented Filter.''}
\emph{Proceedings of the 2003 American Control Conference, 2003.} 3:
2441--2446 vol.3. \url{https://doi.org/10.1109/ACC.2003.1243441}.

\bibitem[\citeproctext]{ref-dermerwe2001sq-ut}
Van der Merwe, R., and E. A. Wan. 2001. {``The Square-Root Unscented
Kalman Filter for State and Parameter-Estimation.''} \emph{2001 IEEE
International Conference on Acoustics, Speech, and Signal Processing.
Proceedings (Cat. No.01CH37221)} 6: 3461--3464 vol.6.
\url{https://doi.org/10.1109/ICASSP.2001.940586}.

\bibitem[\citeproctext]{ref-zhou2026time}
Zhou, Xingyu, Roberto Armellin, Dong Qiao, and Xiangyu Li. 2026.
{``Time-Varying Directional State Transition Tensor for Orbit
Uncertainty Propagation.''} \emph{Journal of Guidance, Control, and
Dynamics}, 1--17.

\end{CSLReferences}




\end{document}
